% ============================================
% 《请把手弄脏：AI时代的学与玩，或一座桥的哲学》
% 作者：ColdCat（大连交通大学）
% 编译器：XeLaTeX
% ============================================

\documentclass[UTF8,a4paper,openany,12pt]{ctexbook}

% ====== 页面布局 ======
\usepackage[top=2.5cm,bottom=2.5cm,left=2.8cm,right=2.8cm]{geometry}
\usepackage{fancyhdr}
\pagestyle{fancy}
\fancyhf{}
\fancyhead[LE,RO]{\thepage}
\fancyhead[RE]{\leftmark}
\fancyhead[LO]{\rightmark}
\renewcommand{\headrulewidth}{0.4pt}
\setlength{\headheight}{15pt}

% ====== 颜色系统 ======
\usepackage[table]{xcolor}
\definecolor{ink}{RGB}{31,36,48}        % 正文墨色
\definecolor{muted}{RGB}{105,113,134}   % 次要文字
\definecolor{blue}{RGB}{84,119,196}     % 链接色
\definecolor{blue-soft}{RGB}{234,241,254} % 浅蓝底
\definecolor{gold}{RGB}{255,244,194}    % 高亮底
\definecolor{line}{RGB}{229,232,240}    % 分割线
\definecolor{paper}{RGB}{252,252,253}   % 页面底色
\definecolor{red-soft}{RGB}{255,240,240} % 警示底

% ====== 图表与插图 ======
\usepackage{graphicx}
\usepackage{subcaption}
\usepackage{tikz}
\usepackage{pgfplots}
\pgfplotsset{compat=1.18}
\usepackage{float}
\usepackage[export]{adjustbox}

% ====== 表格 ======
\usepackage{booktabs}
\usepackage{array}
\usepackage{multirow}
\definecolor{triz-blue}{RGB}{234,241,254}
\definecolor{triz-gold}{RGB}{255,244,194}

% ====== 代码块（避免 shell-escape 依赖） ======
\usepackage{listings}

% ====== 引用与链接 ======
\usepackage{hyperref}
\hypersetup{
    colorlinks=true,
    linkcolor=blue,
    urlcolor=blue,
    citecolor=blue
}
\usepackage[backend=biber,style=ieee,sorting=none]{biblatex}
\addbibresource{bib/refs.bib}

% ====== 自定义环境 ======
\usepackage[many]{tcolorbox}
\tcbuselibrary{skins,breakable}

% 课堂原声
\newtcolorbox[auto counter]{classroom-scene}{
    colback=gray!5,
    colframe=gray!50,
    fontupper=\sffamily\small,
    title={\textbf{课堂原声} \hfill \texttt{\thetcbcounter}},
    breakable,
    before skip=12pt,
    after skip=12pt
}

% 师宇手记（反思旁白）
\newtcolorbox[auto counter]{shiyu-note}{
    colback=blue-soft,
    colframe=blue,
    fontupper=\kaishu\small,
    title={\textbf{师宇手记} \hfill \texttt{\thetcbcounter}},
    breakable,
    before skip=18pt,
    after skip=18pt,
    left=8pt,
    right=8pt
}

% 关键洞察（拍案时刻）
\newtcolorbox[auto counter]{insight}{
    colback=gold,
    colframe=gold!80!black,
    fontupper=\bfseries,
    title={\textbf{关键判断} \hfill \texttt{\thetcbcounter}},
    breakable,
    before skip=16pt,
    after skip=16pt
}

% 技术注释
\newtcolorbox[auto counter]{tech-note}{
    colback=red-soft,
    colframe=red-soft!80!black,
    fontupper=\small,
    title={\textbf{技术说明} \hfill \texttt{\thetcbcounter}},
    breakable
}

% 反向决策（每章末尾）
\newtcolorbox[auto counter]{decision-point}{
    colback=paper,
    colframe=blue!50,
    fontupper=\itshape,
    title={\textbf{教师决策时刻} \hfill \texttt{\thetcbcounter}},
    breakable,
    before skip=20pt,
    after skip=20pt
}

% 双声部：师宇旁白（行内）
\newenvironment{shiyu-inline}
  {\begin{quote}\kaishu\small\color{muted}}
  {\end{quote}}

% ====== 标题格式 ======
\ctexset{
    chapter = {
        format = \LARGE\bfseries\centering,
        name = {第,章},
        number = \arabic{chapter},
        aftername = \hspace{1em},
        beforeskip = 40pt,
        afterskip = 30pt
    },
    section = {
        format = \Large\bfseries,
        beforeskip = 24pt,
        afterskip = 12pt
    },
    subsection = {
        format = \large\bfseries,
        beforeskip = 18pt,
        afterskip = 8pt
    }
}

% ====== 其他 ======
\usepackage{enumitem}
\setlist{nosep,left=2em}
\usepackage{epigraph}
\setlength{\epigraphwidth}{0.7\textwidth}

% ====== 文档信息 ======
\title{请把手弄脏：AI时代的学与玩，或一座桥的哲学}
\author{ColdCat（大连交通大学）}
\date{\today}

\begin{document}

% ====== 封面 ======
\begin{titlepage}
\centering
\vspace*{4cm}

{\Huge\bfseries 请把手弄脏}

\vspace{0.8cm}

{\LARGE AI时代的学与玩，或一座桥的哲学}

\vspace{2cm}

{\large 从一张画错的弯矩图开始}

\vspace{0.8cm}

{\large 写给准备重做一节课的高校教师}

\vspace{3cm}

{\large ColdCat（大连交通大学）}

\vfill

{\small\color{muted} 2026年7月 · 公开修订稿}

\end{titlepage}

% ====== 扉页引言 ======
\thispagestyle{empty}
\vspace*{3cm}

\begin{center}
{\large\itshape
``而且过程中就发现，你就真的得是自己会才行，\\
要不然你也不知道 AI 做得对不对。''}

\vspace{1cm}
——学生项目结题汇报，2026 年 7 月 1 日 00:25:20
\end{center}

% ====== 目录 ======
\tableofcontents

% ====== 阅读说明 ======
\chapter*{阅读说明}
\addcontentsline{toc}{chapter}{阅读说明}

\begin{center}
{\large\itshape 这本书把课堂观察、实践案例、研究证据和作者判断分开书写。}
\end{center}

\vspace{1cm}

\begin{quote}
\begin{itemize}[left=0pt]
  \item \textbf{课堂原声}标注日期和时间戳，但转写仍可能有误，关键引文需要回听复核。
  \item \textbf{实践案例}说明发生了什么，不自动证明学习效果。
  \item \textbf{研究证据}用于说明已有研究支持到哪里，也说明不能外推到哪里。
  \item \textbf{作者判断}可以鲜明，但不会伪装成数据。
\end{itemize}
\end{quote}

\newpage

% ====== 序言 ======
\chapter*{序言：当答案变得便宜}
\addcontentsline{toc}{chapter}{序言：当答案变得便宜}

2026 年 7 月 1 日，一门桥梁工程课程设计课进行结题汇报。学生展示的不再只是计算书和图纸：有桥梁拆除游戏，有影响线动画，有交互式计算工具，也有通过 AI 连接建模软件完成的桥梁模型。

作品很丰富，问题也同样具体。一个小组发现，AI 生成的弯矩图画错了；公式被放在不合适的位置；影响线的负值没有参与计算；继续修改时，一个失败版本还覆盖了接近完成的版本。讲到最后，学生说：

\begin{shiyu-inline}
``而且过程中就发现，你就真的得是自己会才行，要不然你也不知道 AI 做得对不对。像那个弯矩图。''

\hfill ——课程结题汇报，2026 年 7 月 1 日 00:25:20
\end{shiyu-inline}

这句话比“AI 会不会取代教师”更接近本书要处理的问题。

生成式 AI 已经可以迅速给出文本、代码、图片和模型。于是，\textbf{交出一个看起来完整的成品}不再足以证明学生经历了学习。反过来，禁止学生使用 AI 也不能恢复旧秩序：学生毕业后面对的正是开放工具、真实任务、多人协作和不断变化的工作流。

真正需要重做的是课堂里的学习回路。

本书所说的学习回路，是一个朴素的过程：\textbf{学生先尝试，系统或他人给出反馈，学生据此修改，然后解释自己为什么这样改。}在桥梁工程里，作品可以是一座游戏中的桥、一张弯矩图、一个有限元模型、一段施工动画或一份方案比较。作品不是终点，而是让思考变得可见的介质。

我过去的课堂更接近常见的工科教学：知识按章节展开，公式在黑板上推导，学生在作业和考试里复现。这样的教学能传递系统知识，却不保证知识进入工程直觉。2026 年 5 月的一次分享中，我谈到一个不太体面的观察：一些已经学过材料力学、结构力学和理论力学的学生，第一次打开搭桥游戏时，仍主要依靠“加一根杆试试”的方式探索。这个观察不能证明理论课无效。它提醒我：\textbf{会在纸面上求解，不等于会在一个开放问题中识别什么值得算。}

游戏、可视化和 AI 因此进入课堂。它们不是为了让课程显得新，也不是给公式加一层糖衣。它们分别解决三类问题：游戏压缩尝试与反馈之间的时间；可视化让不可见的关系可以操作；AI 降低制作原型的门槛，让更多学生能把想法变成可检查的作品。

这些工具也会制造新的错觉。画面漂亮，不等于模型正确；学生参与热闹，不等于理解发生；AI 能运行一段代码，不等于代码表达了正确的工程假设。关于数字游戏的研究提示，游戏对学习的平均作用为正，但效果取决于具体设计，而不是“用了游戏”这件事本身\cite{clark2016games}。关于反馈的研究也一再指出，反馈的类型、时机和学生是否真正使用反馈，比“给了多少评语”更重要\cite{shute2008feedback,carless2018literacy}。

所以，这不是一本成功案例集。书中会保留没做出来的自动桥型、画错的弯矩图、无法解释的模态、失控的版本和教师并不满意的作业。失败不是为了制造戏剧，而是因为课堂设计只有在这些地方才显出结构。

\section*{这本书写给谁}

第一读者是高校工科教师。你有自己的专业课，知道什么不能讲错；你可能已经试过几个 AI 工具，也可能正被“全面拥抱 AI”和“严禁 AI 代写”两种声音同时拉扯。你需要的不是更多工具名，而是一套能用于下一节课的设计方法。

教师发展工作者、专业负责人和研究生助教也可以把本书当作磨课材料。桥梁工程是贯穿全书的现场，但书中的任务、反馈、证据链和评价方法可以迁移到机械、建筑、交通、物理实验以及其他强调模型与责任的课程。

\section*{读完以后，请做一件事}

不要重做整门课。选一个学生经常“会算但不会判断”的知识点，把它改成一个 45--90 分钟的小任务。要求学生留下初次尝试、得到的反馈、一次修改和一段解释。允许使用 AI，但要求标明它参与了哪一步、学生怎样核验。

如果这个单元能让你看见学生原本藏在最终答案后面的思考，本书就完成了它最重要的工作。

桥在这里仍然是一个好隐喻，但我会谨慎使用它。桥梁把分开的地方建立成关系；课堂设计也把知识、行动、反馈和责任连在一起。隐喻不能替代证据，正如漂亮的渲染不能替代计算。它只负责提醒我们：真正重要的不是把内容从这一岸搬到另一岸，而是让学生能够自己走过去。

\begin{flushright}
ColdCat（大连交通大学）\\
2026 年 7 月，于大连
\end{flushright}


% ====== 第一部分：梁桥 ======
\part[梁桥]{梁桥\\[4pt] {\large\color{muted}最朴素的跨越：AI 如何学习，人如何学习}}
% ============================================
% 第1章：机器学会了搭桥，但它不知道自己在搭桥
% 梁桥 · 最朴素的跨越
% ============================================

\chapter{机器学会了搭桥，但它不知道自己在搭桥}

\vspace{0.5cm}

\begin{center}
{\small\color{muted} 梁桥是最简单的桥。一根梁，两个支座，就跨越了。}
{\small\color{muted} 但最简单的，往往最难做好。}
\end{center}

\vspace{1cm}

\section{AI是怎么学习的？}

2026年7月2日，我在准备一场报告。报告的名字叫“AI时代，教学的变与不变”。我对着录音笔说了一段话，后来被转成了文字：

\begin{shiyu-inline}
“AI的学习方式咱们可以借鉴。它是怎么学习的呢？就是通过训练，给它大量的文本。很多时候它也不理解这东西是啥，或者说我们也很难说明什么是'理解'。在这种情况下，就不能说明什么是理解。那就是输入大量的文本，在里面去进行贝叶斯网络的转换。可能这个以后就是未来的理解，到底什么是理解呢？那么这就是理解。”

\hfill ——2026年7月2日，录音转写
\end{shiyu-inline}

这段话里有一个词，我自己说出来的时候都没意识到它有多重要：\textbf{“它也不理解这东西是啥。”}

AI的学习方式，拆开来看，其实很简单。

\subsection{第一步：喂数据}

你给AI看海量的文本。桥梁工程教材、设计规范、施工方案、论文、博客、论坛帖子\ldots\ldots 它看过的桥，比任何桥梁工程师一辈子见到的都多。

\subsection{第二步：找规律}

AI在这些文本中寻找统计规律。它发现“简支梁”这个词后面经常跟着“跨中弯矩”；它发现“$qL^2/8$”这个公式常常出现在“简支梁”和“弯矩”附近；它发现“悬索桥”和“主缆”“锚碇”“吊杆”这些词有很高的共现概率。

\subsection{第三步：生成输出}

当你问它“简支梁跨中弯矩是多少”，它做的事情不是回忆、不是推导、不是理解——它是\textbf{计算}：在它见过的所有文本里，最有可能的、最合理的、最符合统计规律的输出，就是“$qL^2/8$”。

它不知道$q$是荷载，不知道$L$是跨度，不知道那个$1/8$是从哪里来的。它只知道——这个答案，命中的概率最高。

\subsection{AI的“记忆”}

在同一段录音里，我还说了另一句话：

\begin{shiyu-inline}
“机器人有上下文，我们经常说学生记得这个就忘了那个。用AI的话说，就是上下文比较短，执行力比较差，情绪价值要求比较高，月费好几千。”

\hfill ——2026年7月2日，同一段录音
\end{shiyu-inline}

AI的“记忆”有一个硬边界，叫\textbf{上下文窗口}。你可以把它想象成一个只能装固定数量文字的盒子。盒子满了，旧的被挤出去。AI只能“看见”盒子里的东西。

这个盒子在2025年大约能装15万字，2026年能装100万字。但不管多大，它仍然是一个盒子。AI没有“经验”。它不会说“我记得大一那年，我的结构力学老师说过\ldots\ldots”它不会说“上次我做这个的时候，犯了一个错误，这次我不会再犯。”

\begin{insight}
\label{insight:chap1}

\textbf{AI的学习，本质上是统计推断。}它的目标函数只有一个：让预测误差最小化。

它不在乎“桥”是什么。它不知道塔科马大桥在1940年倒塌的时候，工程师们震惊地发现，他们引以为傲的静力分析在风的作用下完全失效了。它不知道那些失败改变了一整个学科。

AI知道事实。它不知道事实的重量。
\end{insight}

\section{人是怎么学习的？}

2025年11月4日，下午3点25分，我在课堂上讲涡振原理。

\begin{classroom-scene}
\label{scene:chap1-vortex}

{
\noindent \textbf{2025年11月4日，大连交通大学，桥梁抗风课。}

\vspace{0.3cm}

\noindent 我在电子白板上放了一个CFD模拟视频。一个圆柱体，均匀的风吹过来，尾部交替脱落着涡。我指着屏幕上周期性波动的升力曲线，说：

\noindent “这个，就是涡振。”

\noindent 然后我停了一下，说了一段话，后来转写出来是这样的：

\vspace{0.3cm}

\begin{quote}
“其实我不觉得任何这个荷载特性啊，或者是风震，是你们必须学的。我更希望就是你们看到这个研究的背后——其实就是对一个简单的动力学方程，今天怎么处理，怎么去简化。把这种复杂的问题简化到那去，包括怎么去面对这么复杂的这个问题，怎么把它分解成几个独立的振动形式，然后再去讨论\ldots\ldots 所以我觉得这个反而是很很有意义的。就是一个最简单的公式，怎么能 cover 到这么多的这个现象。”
\end{quote}

\vspace{0.3cm}

\noindent ——\href{2025/文字转写/2025-11-04\_1524\_结构工程课程之涡振（涡震）原理与特性讲解.md}{2025年11月4日，涡振原理课，录音转写}
}
\end{classroom-scene}

这段话里，我没有说“今天的教学目标是理解涡振原理”。我说的是：\textbf{怎么把一个复杂问题分解成几个独立的振动形式。}

这不是知识。这是一种\textbf{思维方式}。而思维方式，AI不会教。

\subsection{人学习时，身体在动，手心在出汗}

2026年5月22日，我在面向呼和浩特铁路局教师的线上分享会上，讲游戏化教学。我讲到了一件事，后来被转写成了文字：

\begin{shiyu-inline}
“学生在上这门课之前已经学过材料力学、结构力学、或者理论力学，但是他搭建这个桥梁的过程或者状态，跟他六岁的小孩搭建这个状态是完全一致的。我们不是说他这个多么低龄或多么幼稚。我们在说，其实我们讲授的很多的理论的这种课程，如果没有经过这种实操、试错、迭代、失败成功的这个过程，他可能很难入到你一个学生的这个直觉里面去。”

\hfill ——2026年5月22日，游戏化教学分享，录音转写
\end{shiyu-inline}

这句话里有一个特别重要的词：\textbf{直觉}。

一个学完了材料力学、结构力学、理论力学的大学生，当他第一次打开Poly Bridge搭桥的时候，他的行为和一个六岁小孩\textbf{完全一样}。他不是在调用公式，他是在用直觉——这里加一根杆试试，那里加高一点试试，塌了，换一种方式，又塌了，再换一种\ldots\ldots

这不是退化。这是\textbf{学习的本质状态}。

\begin{insight}
\label{insight:chap1-intuition}

\textbf{人的学习不是模式匹配。人的学习是——经历从不知道到知道的过程，并被这个过程改变。}

AI的输出是“答案”。人的输出是“我懂了”。AI不会在输出正确结果时心跳加速。AI不会在深夜想起自己第一次搭桥成功的那一瞬间，嘴角微微上扬。
\end{insight}

\subsection{失败的权利}

在同一场分享里，我还说了一句话：

\begin{shiyu-inline}
“学生在学习过程中从来没有失败的权利，他也就没有成功的那个喜悦或者荣耀。”

\hfill ——2026年5月22日，同一场分享
\end{shiyu-inline}

这句话，是我在讲完Poly Bridge之后说的。我给学生看了一段视频：一座桥，小车开上去，桥塌了。然后学生改了三个地方，又试了一次。桥站住了。

整个过程不到两分钟。但这两分钟里，学生经历了：\textbf{困惑→尝试→失败→观察→调整→再尝试→成功→骄傲。}

在传统的课堂里，这个回路通常被压缩成一个步骤：\textbf{“记住这个公式，考试会考。”}

\section{两种学习，两种“理解”}

\begin{table}[H]
\centering
\begin{tabular}{p{0.35\textwidth} p{0.35\textwidth}}
\toprule
\textbf{AI的学习} & \textbf{人的学习} \\
\midrule
目标：最小化预测误差 & 目标：理解世界，并能在其中行动 \\
\addlinespace
输入：海量文本 & 输入：有限但多模态的体验（视觉、触觉、情感） \\
\addlinespace
记忆：上下文窗口内的完美复制 & 记忆：选择性的，在遗忘中保留本质 \\
\addlinespace
反馈：损失函数数值下降 & 反馈：\textbf{桥塌了、老师点头了、心里“啊哈”了一声} \\
\addlinespace
错误：统计偏差，需要更多数据 & 错误：\textbf{最有价值的老师} \\
\addlinespace
“理解”：能生成正确的输出 & “理解”：能解释、能迁移、\textbf{能在自己错的时候知道自己错了} \\
\bottomrule
\end{tabular}
\caption{AI的学习 vs 人的学习}
\label{tab:ai-vs-human}
\end{table}

\section{梁桥的隐喻}

梁桥是所有桥型里最简单的。一根梁，两个支座。它不优雅，不壮观，不需要任何高级的力学原理。它只是\textbf{直接跨越}。

但梁桥有一个致命的弱点：\textbf{跨度越大，自重越大。}

2025年10月21日，我在课堂上讲到这件事：

\begin{shiyu-inline}
“我们的桥梁的载重比其实是非常小的。一般来说，一个大跨度桥梁，大概百分之八十是它的自重，百分之二十是车辆。所以说，我们实际上只有百分之二十的承载能力是用在了承载荷载上。所以这个就是很大的浪费嘛。就是姜文的电影是吧——为了这碟醋包了顿饺子。”

\hfill ——2025年10月21日，桥梁课程，录音转写
\end{shiyu-inline}

“为了这碟醋包了顿饺子”——这句话，是我课堂上说的。后来这段录音被转写成文字，我重新读到的时候，发现它比我想象的更好。

\textbf{梁桥的自重，就是AI时代教育面临的问题。}

我们用80\%的精力教学生“知识”——公式、规范、软件操作。这些东西，AI一秒钟就能生成。我们只有20\%的精力用在真正重要的东西上——\textbf{困惑、尝试、失败、骄傲、判断、审美、责任}。

就像梁桥，80\%的自重，20\%的荷载。\textbf{效率极低。}

但这不是AI的错。AI只是把这个问题暴露了出来。它把“做出来”的门槛降到了零，于是我们突然发现：\textbf{原来我们花了80\%的精力教的东西，根本不值钱。}

\begin{insight}
\label{insight:chap1-beam}

\textbf{梁桥的困境，就是当代教育的困境。}

不是AI太强了，是我们一直在教AI擅长的东西。当AI能生成一切“看起来像桥”的东西，真正稀缺的就不再是生成，而是\textbf{判断}。不再是速度，而是\textbf{责任}。不再是看起来像，而是\textbf{真的站得住}。
\end{insight}

\section{这一章，我们站到了哪里}

这本书叫《请把手弄脏》。这五个字来自一个很具体的画面：一个学生，打开电脑，自己写了一个有限元求解器，然后指着它说——“这是我的。”

这个画面出现在这本书的尾声。但在这之前，我们还有很多路要走。

这一章，我们只做了一件事：\textbf{把AI的学习和人的学习，放在一起看。}

AI的学习是统计推断。人的学习是意义生成。AI输出的是“答案”。人输出的是“我懂了”。AI不会骄傲。AI不会在深夜想起自己第一次搭桥成功的那一瞬间。

梁桥是最简单的桥。但最简单的，往往最难做好。

\begin{shiyu-note}
\label{note:chap1}

\textbf{AI已经会搭桥了，问题是人还想不想过桥。}

“AI的学习”这个词，本身就是一个陷阱。我们把“输出正确答案”和“理解”混为一谈了。就像我们把“会背乘法表”和“懂数学”混为一谈一样。

AI时代的教育，真正的问题不是“AI能替代什么”，而是\textbf{“当AI能替代一切，人还剩下什么？”}

人剩下的，恰好是AI永远无法拥有的东西：困惑、好奇心、挫败感、骄傲、责任感、审美判断、以及“非做不可”的理由。

这些，才是学习的真正内容。而AI的出现，只是帮我们把这些东西从“知识传递”的杂音中清晰地分离了出来。
\end{shiyu-note}

\vspace{1cm}

\begin{decision-point}
\label{decision:chap1}

\noindent\textbf{这一章，我们看到了两种“学习”的本质区别。}

\vspace{0.3cm}

\noindent 现在，打开你的下一节课的PPT。找到第一页。删掉“课程目标”。换成一句话，让你的学生\textbf{想}知道接下来会发生什么。

\noindent 不是“通过本章学习，你将掌握\ldots\ldots”
\noindent 而是“你有没有想过，为什么\ldots\ldots？”

\vspace{0.3cm}

\noindent 就这一句。改完发给我。
\end{decision-point}

\vspace{0.5cm}

\begin{flushright}
{\small\color{muted} 下一章：为什么孩子会蹲在沙坑里搭三个小时的桥——关于游戏的本质。}
\end{flushright}
% ============================================
% 第2章：为什么孩子会蹲在沙坑里搭三个小时的桥
% 梁桥 · 游戏本质
% ============================================

\chapter{为什么孩子会蹲在沙坑里搭三个小时的桥}

\vspace{0.5cm}

\begin{center}
{\small\color{muted} 一个学完了材料力学、结构力学、理论力学的大学生，}
{\small\color{muted} 当他第一次打开 Poly Bridge 搭桥的时候，}
{\small\color{muted} 他的行为和一个六岁小孩完全一样。}
\end{center}

\vspace{1cm}

\section{洗浴中心里的桥梁工程师}

2026年5月22日，我在一场面向呼和浩特铁路局教师的线上分享会上，讲了一段话，后来被转写成了文字：

\begin{shiyu-inline}
``学生在洗浴的时候拿个平板，然后一直在玩这个游戏，一直在建不同的桥梁，玩到这个洗浴都关门了。那可能九十点钟还在这个沉浸的这个过程。''

\hfill ——2026年5月22日，游戏化教学分享，录音转写
\end{shiyu-inline}

一个大学生，在洗浴中心，拿着平板，搭桥搭到洗浴关门。

这不是一个段子。这是真实发生的事情。我的学生，在大连的某家洗浴中心，泡完澡之后，没有刷短视频，没有打王者荣耀，而是在用 Poly Bridge 搭桥。他在搭一座斜拉桥。他调了三次缆索的张力，改了两次桥塔的高度，最后一辆小车稳稳地开了过去。他抬头一看，洗浴中心要关门了。

我后来问过他：``你当时在想什么？''

他说：``没想什么。就是塌了，想让它站住。''

\section{``没有人说我用AI去糊弄帮我打游戏''}

在同一场分享里，我还说了另一句话：

\begin{shiyu-inline}
``没有人说我用AI去糊弄帮我打游戏。如果先去下棋或者打牌，没有人说我去用AI来帮我下棋或者打牌。因为我们本身就喜欢这种活动。那我们能不能把这个打牌呀、打游戏呀这种状态，让他们在课堂上也有这种状态，然后就是沉浸在这里面，不会去想我怎么去糊弄他，而是真正的喜欢他。''

\hfill ——2026年5月22日，同一场分享
\end{shiyu-inline}

这句话里有一个特别重要的对比：\textbf{学生会用AI糊弄一份作业，但没有人会用AI糊弄一场游戏。}

为什么？因为作业是别人让你做的事。游戏是你自己选择做的事。作业的目标是``完成''。游戏的目标是``赢''——或者说，是``经历从不会到会的过程，并且在这个过程中感受到自己正在变强''。

这不是我的理论。这是任何一个在洗浴中心搭桥搭到关门的大学生，用他的身体告诉我的。

\section{``很热闹，但是学生没学到什么''}

但游戏化不是没有争议的。

2026年5月31日，我做了另一场游戏化教学分享。开场前，我用一个``六顶思考帽''的活动，收集了老师们对游戏化的看法。后来这些话被转写成了文字：

\begin{shiyu-inline}
``游戏化有哪些价值、好处和收益呢？提高学生积极性啊，形式上比单纯讲授新颖一点，体验性强一点，兴趣是最好的老师，吸引学生的兴趣，提升课堂质量，让这个孩子做课堂的主人。''

``消极的部分——课堂纪律难以把控。学习重点可能会偏移。很热闹，但是学生没学到什么。玩到位了，是不是学到位了？还有就是，领导可能会不喜欢，影响考核。学生太活跃，领导不喜欢，影响我考核，降低我薪资。''

\hfill ——2026年5月31日，游戏化教学设计分享，录音转写
\end{shiyu-inline}

这些话，我每次分享都会听到。它们分成两派：

\textbf{积极派说：}游戏化能提高兴趣、增加参与、让学生做课堂的主人。

\textbf{消极派说：}游戏化会让课堂失控、目标偏移、热闹但没学到东西、领导不喜欢。

两派说的都对。但两派都漏掉了一个问题：\textbf{游戏为什么好玩？}

\section{游戏的核心：即时反馈}

2026年5月22日，我在讲 Poly Bridge 的时候，说了一段话：

\begin{shiyu-inline}
``其实我们非常关注的就是游戏它本身有什么特征，跟我们传统的教学有什么本质的区别。其实我感觉最重要的就是\textbf{及时反馈}。就是我做出来一个操作，马上我就能获得这个东西行还是不行，好还是不好，这样的一个回应。然后我跟这个回应，再去做这个调试。''

\hfill ——2026年5月22日，游戏化教学分享
\end{shiyu-inline}

这句话是我对游戏本质的理解：\textbf{即时反馈。}

你动一下，桥马上有反应。你加一根杆，桥的颜色变了。你放一辆小车，桥塌了——不是考试之后才知道，是\textbf{当场就知道}。你改了三个地方，又放一辆小车，桥站住了。你心里``啊''了一声。

这个回路——\textbf{动作→反馈→调整→再动作}——在游戏里，几秒钟就能完成一轮。在传统的课堂里，这个回路可能需要几周。你听了一节课，做了一次作业，交上去，等老师批改，拿回来，看到分数——但你已经忘了你当时是怎么想的了。

\subsection{失败的权利}

在同一场分享里，我还说了一句后来被反复引用的话：

\begin{shiyu-inline}
``学生在学习过程中从来没有失败的权利，他也就没有成功的那个喜悦或者荣耀。''

\hfill ——2026年5月22日，同一场分享
\end{shiyu-inline}

这句话是我在展示一段视频时说的。视频里，一座桥塌了。然后学生改了三个地方，桥站住了。整个过程不到两分钟。

在传统的课堂里，失败是\textbf{被惩罚的}。你答错了，扣分。你作业做错了，扣分。你考试没考好，扣分。但在游戏里，失败是\textbf{免费的}。你塌了一百次桥，没有人扣你的分。你只需要再试一次。每一次失败都告诉你一件事：\textbf{这里不对，换一种方式。}

\begin{insight}
\label{insight:chap2}

\textbf{游戏不是逃避。游戏是主动选择困难，并在安全的环境中反复失败，直到成功。}

孩子在沙坑里搭三个小时的桥，不是因为好玩——是因为他正在经历一个他主动选择的困难，并且他正在克服它。每一次沙堡塌了，他都知道为什么塌。每一次沙堡站住了，他都感到一种发自心底的骄傲。

这种骄傲，考试得100分给不了。
\end{insight}

\section{``我们从来没有教过学生从零开始跟工程师一起去思考''}

2026年5月31日，我讲了另一段话：

\begin{shiyu-inline}
``我们传统的教学里面从来没有教过学生从零开始跟工程师一起去思考这个问题。之前我们就告诉说，应该是这么做的是吧？这个桁架一般这么高，五米高，二十米的桥应该是五米高的桁架。这个构件应该什么尺寸。这个状态是很可悲的——因为他在设计别人做的结构。他从来没有真正像一个工程师一样，说我们思考一下这个问题该怎么解决。它可能不一定只有一个解法。''

\hfill ——2026年5月31日，游戏化教学设计分享
\end{shiyu-inline}

这段话里有一个词，我自己说出来的时候都觉得有点重：\textbf{可悲}。

我们教学生``应该怎么做''。我们告诉他们桁架应该多高，梁应该多宽，配筋应该怎么配。我们从来没有让他们自己想过：\textbf{如果让你来设计，你会怎么做？}

但在游戏里，学生第一次被放在了\textbf{设计者}的位置上。不是``记住这个公式，考试会考''，而是``这里有一道沟，你有一座桥要搭，成本不能超过三万，车要能过，你自己想办法''。

这个区别，就是\textbf{学习}和\textbf{被灌输}之间的区别。

\section{``六岁小孩''和``大学生''}

2026年5月22日，我说了那句话：

\begin{shiyu-inline}
``学生在上这门课之前已经学过材料力学、结构力学、或者理论力学，但是他搭建这个桥梁的过程或者状态，跟他六岁的小孩搭建这个状态是完全一致的。我们不是说他这个多么低龄或多么幼稚。我们在说，其实我们讲授的很多的理论的这种课程，如果没有经过这种实操、试错、迭代、失败成功的这个过程，他可能很难入到你一个学生的这个直觉里面去。''

\hfill ——2026年5月22日，游戏化教学分享
\end{shiyu-inline}

这句话里有一个非常重要的词：\textbf{直觉}。

一个学完了材料力学、结构力学、理论力学的大学生，当他第一次搭桥的时候，他的行为和一个六岁小孩\textbf{完全一样}。他不是在调用公式，他是在用直觉——这里加一根杆试试，那里加高一点试试，塌了，换一种方式。

这不是退化。这是\textbf{学习的起点}。所有的理论，所有的公式，所有的规范，都必须经过这个起点——\textbf{动手做，失败，调整，再动手}——才能真正进入一个人的直觉。没有这个过程，你背再多的公式，那个公式也只是黑板上的一行字，不是你手指上的感觉。

\section{梁桥的隐喻之二}

梁桥是最简单的桥。但最简单的，最难做好。因为梁桥的跨度越大，自重越大——80\%的自重，20\%的荷载。

\textbf{游戏的本质，就是帮我们解决梁桥的困境。}

游戏不是``加糖''。游戏不是``让学习变轻松''。游戏是\textbf{重新设计学习回路}——把困惑→尝试→失败→调整→成功的循环，从几周压缩到几分钟。让学生在安全的环境中反复失败，直到他们自己发现那个公式的含义。

当一个学生蹲在沙坑里搭三个小时的桥，他不是在``玩''。他是在\textbf{学习}——用人类最古老、最自然、最有效的方式学习。

\begin{shiyu-note}
\label{note:chap2}

\textbf{游戏不是学习的敌人。游戏是学习最原始的形态。}

我们总是把``学习''和``严肃''捆绑在一起。好像不够严肃的东西，就不可能是学习。但所有三岁孩子都知道：学习是一件让人上瘾的事。你学走路，摔倒了，爬起来，再摔，再爬起来——你从来没有觉得``学走路好累，我不学了''。因为走路不是别人让你学的，是\textbf{你自己想学的}。

游戏化教学，不是把课堂变成游乐场。是把课堂重新设计成\textbf{一个可以安全失败、即时反馈、主动探索的地方}。就像沙坑。就像 Poly Bridge。就像任何一个孩子蹲在那里搭了三个小时的桥，然后抬头说：``站住了。''

\end{shiyu-note}

\vspace{1cm}

\begin{decision-point}
\label{decision:chap2}

\noindent\textbf{这一章，我们看到了游戏的本质：即时反馈、安全失败、主动选择困难。}

\vspace{0.3cm}

\noindent 现在，回顾你上一次批改的作业。你花了多少时间批改？学生拿到批改结果之后，距离他做作业的时候，已经过了多久？

\noindent 如果这个反馈回路太长，试着缩短它。下一节课，不要等到课后批改。在课堂上，让学生当场做，当场看结果，当场改。就一次。

\noindent 告诉我，发生了什么。
\end{decision-point}

\vspace{0.5cm}

\begin{flushright}
{\small\color{muted} 下一章：当学生带着ChatGPT的答案走进教室——旧系统是如何被冲击的。}
\end{flushright}
% ============================================
% 第3章："这个学了有什么用？"——效率如何杀死了意义
% 梁桥 · 效率陷阱
% ============================================

\chapter{``这个学了有什么用？''——效率如何杀死了意义}

\vspace{0.5cm}

\begin{center}
{\small\color{muted} 梁桥是最简单的。但梁桥有一个致命的缺陷：}
{\small\color{muted} 跨度越大，自重越大。80\%的自重，20\%的荷载。}
{\small\color{muted} 效率极高。意义极低。}
\end{center}

\vspace{1cm}

\section{那个问题}

每一个教土木的老师，都被问过这句话。

可能在课堂上，一个坐在后排的学生，没有举手，直接问了出来。可能在办公室答疑，学生犹豫了很久，最后小声说了一句。可能在毕业聚餐上，几杯酒之后，终于说了实话。

\begin{center}
{\large\itshape ``老师，学这个，到底有什么用？''}
\end{center}

这个问题，我听过很多次。每一次，我都回答。我说，桥梁工程是基础设施的基石。我说，中国有超过一百万座公路桥，每年还需要新建几万座。我说，毕业生的就业率不低，薪资也不差。

但我知道，这些回答，没有一个真正回答了那个问题。

因为学生问的不是``学这个能找到什么工作''。学生问的是：\textbf{在AI能做一切的时代，我花四年时间学一门手艺，这件事本身，还值得吗？}

\section{``为了这碟醋包了顿饺子''}

2025年10月21日，我在课堂上讲桥梁的载重比。我讲了一段话，后来被转写成了文字：

\begin{shiyu-inline}
``我们的桥梁的载重比其实是非常小的。一般来说，一个大跨度桥梁，大概百分之八十是它的自重，百分之二十是车辆。所以说，我们实际上只有百分之二十的承载能力是用在了承载荷载上。所以这个就是很大的浪费嘛。就是姜文的电影是吧——\textbf{为了这碟醋包了顿饺子}。''

\hfill ——2025年10月21日，桥梁课程，录音转写
\end{shiyu-inline}

我当时说这句话的时候，教室里笑了一下。姜文的电影梗，大家都懂。但后来我重新读到这段转写，发现这个比喻比我想象的更深。

\textbf{梁桥80\%的自重，就是当代教育80\%的``无用功''。}

我们常常把大量时间用于公式复述、规范检索、软件操作与考试组织。AI正在加速其中一部分工作，于是那个旧问题变得更响：省下来的时间，能否真正还给尝试、解释、核验和创造？这里不再套用一个没有调查依据的百分比。

我们只有20\%的精力用在真正重要的东西上——\textbf{困惑、尝试、失败、骄傲、判断、审美、责任}。

就像梁桥，80\%的自重，20\%的荷载。\textbf{为了这碟醋，包了顿饺子。}

\section{载重比：一个工程概念如何变成教育隐喻}

在那堂课上，我还讲了另一个故事——关于结构设计大赛：

\begin{shiyu-inline}
``我们做了一个100克的船，它可以承载19公斤的重量。大概载重比是190倍。但是我们的桥梁，大跨度桥梁，载重比是非常小的。一百米的桥，大概百分之八十是它的自重。''

\hfill ——2025年10月21日，同一堂课
\end{shiyu-inline}

\textbf{载重比}，是一个工程概念。它衡量的是：一座桥的重量里，有多少是``必须的负担''（自重），有多少是``真正的价值''（承载交通）。

100克的船，承载19公斤——载重比190倍。100米的桥，80\%是自重——载重比1.25倍。

\textbf{载重比，也是一个教育概念。}

一门课的载重比是多少？学生花了80个小时，有多少小时是在``真正理解''，有多少小时是在``完成要求''？老师花了80个小时备课，有多少小时是在``设计学习体验''，有多少小时是在``做PPT排版''？

梁桥的困境，就是教育的困境：\textbf{我们花了太多精力做那些``必须做但没什么用''的事，以至于没时间做``真正重要''的事。}

\section{AI让这个问题变得更尖锐了}

2026年6月，我在知识星球上看到一段话，是师宇写的：

\begin{shiyu-inline}
``AI时代的教育，不该再假装学生生活在校园平台里。很多教育AI产品喜欢做封闭知识库、智能问答、自测系统。它们假设学生愿意在学校提供的环境里耐心学习。但真实世界不是这样运行的。学生未来面对的是开放工具、真实任务、复杂工作流和不断变化的平台。''

\hfill ——师宇博文，2026年6月8日
\end{shiyu-inline}

AI把这个问题逼到了极致。因为AI不只让``知识传递''变快了——它让``知识传递''\textbf{变得几乎免费}。

你不需要花四年时间学``桥梁工程基础知识''。你只需要打开ChatGPT，输入``帮我设计一座20米简支梁桥''，它就会给你一份完整的计算书。格式规范，内容合理，甚至附带了参考文献——虽然那些参考文献可能有一半是它编的。

\textbf{当AI能把80\%的``自重''瞬间完成，那剩下的20\%——那些AI做不了的东西——到底是什么？}

\section{效率杀不死意义，但它会让意义变哑}

2026年6月11日，师宇写了一篇博文，标题叫``有什么用呢''。他研究了AI视频生成工具——Coze、Seedance、花生AI——兴奋地做了很多尝试。然后被一句话问住了：

\begin{shiyu-inline}
``我原本兴奋于'是不是该认真做一些有趣的视频'，却被一句'有什么用呢'问住。什么算有用？必须产出、变现、转化才叫有用吗？不是所有有用的东西，都必须马上证明自己有用。有些价值正是在暂时无用时显现。''

\hfill ——师宇博文，2026年6月11日
\end{shiyu-inline}

这句话击中了问题的核心。

``这个学了有什么用''——这句话本身，就是效率思维的产物。在效率至上的世界里，一切``无用''之事都被驱逐了。诗、音乐、游戏、搭一座你很确定不会有人走的桥。这些事都没有``用''。但它们让生活值得过。

\textbf{效率杀不死意义。但效率会让意义变哑。}当所有人都只问``有什么用''，那些``没用但重要''的东西就不再被说出口。它们没有消失，它们只是沉默了。

\section{从梁桥到拱桥}

梁桥的困境，本质上是\textbf{效率与意义的矛盾}。

你想让梁桥的跨度更大？你必须让它更重。你想让它更轻？它的跨度就会变小。这是一个物理定律，没有变通的办法。

但工程师没有放弃。他们发明了拱桥。

拱桥的秘密是：\textbf{把垂直的荷载，转化为水平的推力。}压力不是被抵抗的，而是被利用的。石头不是被压垮的，而是被压紧的。\textbf{压力，创造了强度。}

这就是第二部分要讲的事。

在接下来的几章里，我们会看到：当AI冲击旧的教育系统，当学生带着ChatGPT的答案走进教室，当排名、积分、点赞成为新的课堂货币——这些压力，不是要摧毁教育的敌人。它们是\textbf{建造拱桥的石头}。

\begin{insight}
\label{insight:chap3}

\textbf{梁桥的困境，不是要我们放弃梁桥。是要我们学会造拱桥。}

效率杀不死意义。但效率会让我们忘记，意义从来不是``有用''的东西。意义是你在洗浴中心搭桥搭到关门，别人问你为什么，你只能说：``没想什么。就是塌了，想让它站住。''

那就是意义。那不需要被证明有用。
\end{insight}

\vspace{1cm}

\begin{shiyu-note}
\label{note:chap3}

\textbf{在效率的海洋里，为``无用的认真''保留一片海滩。}

有一个问题，我建议所有老师在开学第一节课问学生：

``如果这门课不算学分，没有成绩，不写入档案——你还会来上吗？''

这不是一个测试。这是一面镜子。它照出的不是学生是否勤奋，而是这门课是否值得被教。

如果答案是``不会''，那问题不在学生。在课程设计。我们花了太多时间让课程``有用''——符合大纲、覆盖考点、支撑后续课程——却忘了让课程``值得''。

值得，不是有用。值得是：即使没有任何外部奖励，你仍然愿意做这件事。就像孩子蹲在沙坑里搭三个小时的桥。就像洗浴中心里那个搭桥搭到关门的大学生。

AI时代，效率不再是稀缺品。意义才是。
\end{shiyu-note}

\vspace{1cm}

\begin{decision-point}
\label{decision:chap3}

\noindent\textbf{这一章，我们看到了效率如何杀死了意义。}

\vspace{0.3cm}

\noindent 下一次备课，做一件事：把你准备讲的内容分成两列。左边写``这个学生必须知道''，右边写``这个学生不学也可以，但我想让他们感受一下''。

\noindent 然后，把左边砍掉一半。把右边加倍。

\noindent 告诉我，那节课发生了什么。
\end{decision-point}

\vspace{0.5cm}

\begin{flushright}
{\small\color{muted} 第一部分结束。下一部分：拱桥——当AI冲击旧系统，压力如何创造强度。}
\end{flushright}


% ====== 第二部分：拱桥 ======
\part[拱桥]{拱桥\\[4pt] {\large\color{muted}压力创造强度：新工具如何撞击旧课堂}}
% ============================================
% 第4章：当学生带着ChatGPT的答案走进教室
% 拱桥 · 旧系统坍塌
% ============================================

\chapter{当学生带着ChatGPT的答案走进教室}

\vspace{0.5cm}

\begin{center}
{\small\color{muted} 拱桥的秘密：把垂直的荷载，转化为水平的推力。}
{\small\color{muted} 压力不是被抵抗的，而是被利用的。}
{\small\color{muted} 石头不是被压垮的，而是被压紧的。}
\end{center}

\vspace{1cm}

\section{那一刻}

2025年10月24日，下午5点59分，桥梁风工程课。

我布置了一个课堂练习：用AI分析一段风洞数据，提取前三阶模态频率。三十秒后，一个学生举手。

\begin{shiyu-inline}
``老师，出来了。一阶0.23，二阶0.58，三阶1.04。''

我问：``这个结果，你们觉得对吗？''

教室里安静了五秒。

``呃\ldots\ldots AI 给的，应该对吧？''

\hfill ——2025年10月24日，桥梁风工程课，录音转写
\end{shiyu-inline}

那一刻，我意识到一件事：\textbf{旧系统已经坍塌了。}

旧系统是什么？旧系统是：老师拥有知识，学生需要知识，课堂是知识传递的场所。老师讲，学生记。老师考，学生答。这个系统运行了几百年，所有人都觉得理所当然。

然后ChatGPT来了。

现在，任何一个学生都可以在三十秒内得到一份完整的模态分析结果。格式规范，步骤清晰，甚至附带了解释。如果知识传递是教学的全部内容，那我已经没有存在的必要了。

\section{讲解变便宜了，什么变贵了？}

2026年5月7日，我在一个人工智能与高等教育融合的工作坊上做分享。那场分享持续了三个小时，后来被转写成了文字。中间有一段，我讲到了这个变化：

\begin{shiyu-inline}
``视频和解释随处可得，但让学习者判断'我到底懂没懂'、完成测试、调用知识解决具体问题，才是未来真正值钱的教育体验。这比'AI帮老师备课'深一层——它重新定义了课程的稀缺品。''

\hfill ——2026年5月7日，AI与高等教育融合工作坊，录音转写
\end{shiyu-inline}

这段话里有一个特别重要的词：\textbf{稀缺品}。

在旧系统里，稀缺品是\textbf{知识}。老师拥有知识，学生没有。所以老师站在讲台上，学生坐在下面。这个权力结构建立在信息的\textbf{不对称}上。

AI把这种不对称炸掉了。现在，任何人有手机就能获取任何知识。知识不再是稀缺品。\textbf{判断力才是。}

\begin{insight}
\label{insight:chap4}

\textbf{AI把``讲解''变便宜了。但把``判断''变贵了。}

过去，一个老师最值钱的能力是``讲得清楚''。未来，一个老师最值钱的能力是``设计得出让学生不得不判断的场景''。

不是``我来告诉你答案''。而是``我来制造一个你必须自己判断的情境''。
\end{insight}

\section{旧系统坍塌的三个瞬间}

旧系统不是在一天之内坍塌的。它是在三个瞬间，一点一点瓦解的。

\subsection{第一个瞬间：学生交上来的作业，AI写的}

2025年秋天，我开始在学生的作业里发现一种奇怪的``完美''。格式规范，逻辑清晰，参考文献格式正确——但所有参考文献都是编的。

我找了一个学生来问。他说：``老师，我用AI写的。但是我不知道它编了参考文献。我以为它查过。''

他不是在偷懒。他是真的不知道AI会编造参考文献。在他的认知里，AI就像一个超级搜索引擎——它给你的东西，应该是有出处的。

\textbf{旧系统的第一个裂缝：学生不再区分``知识''和``AI生成的内容''。}

\subsection{第二个瞬间：学生用AI写的作业，比老师写的教案还好}

2025年冬天，我试着用AI生成了一份涡振原理的讲义。我把它和我的教案放在一起比较。AI的版本更简洁，例子更丰富，甚至还配了流程图。

我在课堂上说了一句话，后来被转写成了文字：

\begin{shiyu-inline}
``还凑合啊，这个挺挺像那个像像那回事的啊。咱们一会儿会大概跟大家一起过一下，窝阵的这种振动现象的特点。''

\hfill ——2025年11月4日，涡振原理课，录音转写
\end{shiyu-inline}

我说``还凑合''。但心里想的是：\textbf{它比我写得好。}

\textbf{旧系统的第二个裂缝：老师发现自己的核心能力——``讲清楚''——被AI超越了。}

\subsection{第三个瞬间：学生问了一个AI回答不了的问题}

2026年春天，一个学生在课堂上问了我一个问题。不是``简支梁跨中弯矩是多少''——那个问题AI能回答。他问的是：

``老师，AI说这座桥的梁高应该取1.5米。但我用游戏试了一下，1.2米也能站住。为什么规范说1.5米？''

这个问题，AI回答不了。因为AI不知道规范里的安全系数是怎么来的。它不知道那些系数背后是几十年的工程事故和血的教训。它不知道1.5米不是``最优解''，而是``最安全的解''。

\textbf{旧系统的第三个裂缝：学生问出了一个只有人才能回答的问题。}

\section{拱桥的隐喻：压力如何变成强度}

拱桥的秘密，不是抵抗压力，而是利用压力。

你把一块石头放在拱顶上，它不往下掉——它被两边的石头挤住了。压力越大，挤得越紧。石头不是被压垮的，而是被压紧的。\textbf{压力，创造了强度。}

旧系统坍塌带来的压力，不是要摧毁教育的敌人。它是建造拱桥的石头。

\begin{itemize}
  \item \textbf{AI能生成一切→压力→教师不再靠``讲清楚''吃饭→必须设计``让学生判断''的场景。}
  \item \textbf{学生能随时获取答案→压力→课堂不再是知识传递的场所→必须变成困惑制造的场所。}
  \item \textbf{作业可以用AI完成→压力→作业不再是检验知识的手段→必须变成检验判断的手段。}
\end{itemize}

每一块压力，都对应着一块新的能力。这就是拱桥的建造方式。

\section{教师还剩什么？}

2026年7月2日，我在准备一场报告——``AI时代，教学的变与不变''。我对着录音笔说了一段话：

\begin{shiyu-inline}
``我们怎么去重新评估人的价值呢？其实我觉得比较有意思的，首先就是AI的学习方式咱们可以借鉴。它是怎么学习的呢？就是通过训练，给它大量的文本。很多时候它也不理解这东西是啥。''

\hfill ——2026年7月2日，录音转写
\end{shiyu-inline}

这段话里，我其实是在问自己：\textbf{如果AI不理解，而人理解——那``理解''到底是什么？}

后来我慢慢想明白了。理解不是``能输出正确答案''。理解是：

\begin{itemize}
  \item \textbf{能在自己错的时候，知道自己错了。}
  \item \textbf{能在一个新的情境里，判断旧的知识是否适用。}
  \item \textbf{能解释为什么选这个方案，而不是那个方案。}
  \item \textbf{能指着一样东西说：``这是我的。''}
\end{itemize}

这些，AI全做不到。不是因为它不够聪明。是因为它没有\textbf{身体}。它没有碰过混凝土，没有感受过钢筋被弯折时的阻力，没有在深夜因为一个参数调不对而摔过键盘。它没有\textbf{做过}任何事。它只是\textbf{算过}。

\textbf{教师剩下的，就是那些AI算不出来的东西。}

\begin{shiyu-note}
\label{note:chap4}

\textbf{旧系统坍塌了。不是坏事。}

旧系统本来就有问题。它假设知识的稀缺性，但知识从来都不稀缺——好奇心才是。它假设老师是知识的唯一来源，但老师从来都不是——经验才是。它假设学习是``传递''，但学习从来都不是——学习是\textbf{转化}。

AI只是把这些旧假设炸掉了。它把我们从``知识搬运工''的角色里解放了出来。现在，我们可以做真正重要的事了：\textbf{设计困惑。制造反馈。守护判断。见证骄傲。}

这些事，AI做不了。不是因为它们太难。是因为它们太\textbf{人}了。
\end{shiyu-note}

\vspace{1cm}

\begin{decision-point}
\label{decision:chap4}

\noindent\textbf{这一章，我们看到了旧系统如何坍塌，以及坍塌之后教师还剩什么。}

\vspace{0.3cm}

\noindent 下一节课，做一件事：在你讲完一个知识点之后，不要问``大家听懂了吗''。问：``AI刚才给我的答案，你们觉得对吗？为什么？''

\noindent 你会发现，学生沉默的时间变长了。然后，有人会开口。那个开口的人，说的不是``标准答案''。他说的是\textbf{判断}。
\end{decision-point}

\vspace{0.5cm}

\begin{flushright}
{\small\color{muted} 下一章：排名、积分、点赞——我们在设计``微型制度''。}
\end{flushright}
% ============================================
% 第5章：排名、积分、点赞：我们在设计"微型制度"
% 拱桥 · 激励设计
% ============================================

\chapter{排名、积分、点赞：我们在设计``微型制度''}

\vspace{0.5cm}

\begin{center}
{\small\color{muted} 拱桥的每一块石头都被两边的石头挤住。}
{\small\color{muted} 如果有一块石头想偷懒——它就会掉下来。}
{\small\color{muted} 制度，就是那两边的石头。}
\end{center}

\vspace{1cm}

\section{``分已经打完了，而且都是最高分''}

2026年6月23日，我在跟两位老师讨论我开发的AI教学平台。我讲了一段话，后来被转写成了文字：

\begin{shiyu-inline}
``以前我们在学习通里评分的时候，台上学生往台上一站，你不是可以评分了吗？分已经打完了，而且都是最高分。为啥他打多少分跟自己没关系？你再有一个是你只需要打个分嘛，写个分最省事，人就最小用力原则嘛。那只要打个分，打个满分，啪提交就完事了。''

\hfill ——2026年6月23日，AI教学平台讨论，录音转写
\end{shiyu-inline}

这段话描述的场景，每一个用过学习通评分功能的老师都见过。一个学生上台汇报，台下三十个学生，三十秒内全部打完分。清一色的满分。没有人真的在``评价''。所有人都在``完成''。

为什么？因为\textbf{打多少分，跟自己没关系。}

这就是旧系统里评价机制的致命缺陷：评价者和被评价者之间，没有利益关联。我给你打满分，不花我一分钱。我给你打低分，你也不会知道是我打的。那为什么不打满分呢？最小用力原则。

\section{``他必须打这个分很公正''}

然后我讲了我怎么改的：

\begin{shiyu-inline}
``但是现在我要做的是，因为你给别人打完分之后，别人马上加到积分里，这个积分一变就影响到了你自己的排名，影响了你期末的分数。所以他必须打这个分很公正。''

\hfill ——2026年6月23日，同一场讨论
\end{shiyu-inline}

我改了一个东西：\textbf{把评价者的利益和被评价者的表现绑在一起。}

你给别人打的分数，会影响别人的排名。别人的排名，会影响你的排名。如果你给一个很烂的作品打满分，那个人的排名会上升，你的排名会下降。如果你给一个很好的作品打低分，那个人的排名会下降——但你的排名不会因此上升，因为AI会检测你的评分是不是偏离了均值太远。

\textbf{这不是一个评分系统。这是一个微型制度。}

在这个制度里，``公正''不是一种道德要求，而是一种\textbf{利益计算}。你不需要做一个好人。你只需要做一个理性的人。而理性的人，会打公正的分。

\section{``AI给整理的所有学生对他的评价''}

制度里还有另一个细节：

\begin{shiyu-inline}
``30个评审团60句话。别人看到不是60句话，而是匿名的由AI整理出来的一段。AI给整理的所有学生对他的评价，整理完之后你看不出是哪位学生给他的这个东西，但一定里边你所说的话被采纳进去了。''

\hfill ——2026年6月23日，同一场讨论
\end{shiyu-inline}

我做了两件事：\textbf{匿名化}和\textbf{AI聚合}。

匿名化，让学生敢说真话。没有人知道那句批评是谁写的，所以没有人会因为批评同学而被记恨。

AI聚合，让学生的话被``听见''。三十个人写的六十句话，对一个人来说太多了。但AI把它们整理成一段流畅的反馈，那个学生就能在三十秒内读完，然后马上知道自己的问题在哪，下一版怎么改。

\textbf{制度的目的是让反馈变得可操作。}

\section{``他排名啪一杆就冲上去了''}

然后我讲了另一个设计：抢答机制。

\begin{shiyu-inline}
``第一个答完的学生能得十分，第二个答完能得九分，第三个答完得八分。咱以前课堂上最头疼的是好学生答完之后把答案啪一散，全班都知道。但是你现在没问题，你为了能够快速答完，好学生根本没工夫去管其他人。他赶紧先把自己答案提交了。''

``我上课经常有那样的学生，就是因为这一轮他答的最快最好，直接他就排名啪一杆就冲上去了。就这样，让有可能性他能翻盘的。所以他才愿意参与。要不然他发现自己反正我也超不过前面了，无所谓了，我就是躺平，我不跟你参与了。''

\hfill ——2026年6月23日，同一场讨论
\end{shiyu-inline}

这里有三个设计原则：

\textbf{第一，速度本身就是奖励。}第一个答对的人拿10分，第二个拿9分，第三个8分——后面的人拿到的分越来越低。这个设计把``快''变成了一种优势。好学生不会再把答案分享给全班——因为分享会降低自己的排名。

\textbf{第二，翻盘必须可能。}如果排名永远不变，排在后面的人就会放弃。但一轮抢答就能让排名``啪一杆就冲上去''——这个可能性，是让所有人持续参与的动力。

\textbf{第三，参与本身就有分。}不是只有答对才得分。投票、开放式问题、甚至点``不懂''——只要参与了，AI就给你分。因为\textbf{参与本身就是价值}。一个沉默的课堂，比一个答错的课堂更可怕。

\begin{insight}
\label{insight:chap5}

\textbf{积分系统不是在管理学生。积分系统是在设计课堂的微观制度。}

每一项规则——谁得分、怎么得分、得分之后会发生什么——都在塑造学生的行为。如果你奖励``正确''，学生就会追求安全。如果你奖励``快速''，学生就会追求效率。如果你奖励``参与''，学生就会追求表达。

制度不是中性的。制度本身就在教东西。
\end{insight}

\section{制度的双面性}

2026年6月，在教学理念综合分析中，有一段话让我反复读：

\begin{shiyu-inline}
``积分系统实际上是在设计课堂的'微型制度'。有的平台把学生给他人的评分与自己的排名、期末成绩关联，解决'人人随手满分'；点赞、评论、与教师判断一致还会产生不同积分。这非常聪明，因为它改变了参与的激励结构；也非常危险，因为学生可能开始优化积分、迎合教师或交换人情，而非追求真实判断。''

\hfill ——教学理念综合分析，2026年7月
\end{shiyu-inline}

\textbf{制度是双面的。}

它能让沉默的课堂变活跃，也能让活跃的课堂变表演。它能让学生从不参与到参与，也能让学生从``参与''到``迎合''。它能解决``人人随手满分''，也可能制造``人人优化积分''。

这不是制度的错。任何制度都有双面性。问题不是``要不要制度''，而是\textbf{制度有没有被有意识地设计，以及有没有被持续地审视}。

在我的课上，我保留了那个反思：\textbf{即使平台高度智能化，``人与人的东西''仍被认为最重要。}积分是辅助，不是目的。学生迭代的过程应该由学生自己总结，而不是外包给AI。

\section{拱桥的隐喻之二}

拱桥的每一块石头，都被两边的石头挤住。如果有一块石头想偷懒——它就会掉下来。制度，就是那两边的石头。

但制度不是越多越好。石头太密，拱桥就变成了实心的墙——没有空间，没有呼吸，没有光。制度太松，拱桥就散了——每一块石头都只为自己，整座桥就塌了。

\textbf{好的制度，是让每个人在追求自身利益的同时，恰好做了对整体有利的事。}

就像拱桥。每一块石头都在承受自己的重量。但因为它被两边的石头挤住了，它承受的重量，恰好帮助了整座桥站住。

\begin{shiyu-note}
\label{note:chap5}

\textbf{制度是双面的。但逃避制度的设计，是更危险的。}

很多老师不愿意设计积分系统。他们觉得``学习应该是自发的''，``积分会让学习变味''。但问题是：\textbf{不做设计的课堂，也有制度。}它只是隐性的——学生知道谁在认真听、谁在刷手机、谁会被老师喜欢、谁会被忽视。这些隐性的制度，比显性的积分系统更不公平，因为没有人能质疑它。

把制度显性化，才是公平的开始。你可以看到规则，你可以质疑规则，你可以建议修改规则。制度变成了一件可以被讨论和迭代的东西，而不是一件``从来就是这样''的天经地义。
\end{shiyu-note}

\vspace{1cm}

\begin{decision-point}
\label{decision:chap5}

\noindent\textbf{这一章，我们看到了积分系统如何设计课堂的微型制度。}

\vspace{0.3cm}

\noindent 下一次你让学生互评的时候，不要只让他们打分。加一句话：``你认为这个作品最大的优点是什么？如果你来改，你会改哪里？''

\noindent 然后，把这些评语匿名汇总，用AI整理成一段话，发给那个学生。告诉他：这是三十个人帮你想的改进方向。

\noindent 看看他下一版会改成什么样。
\end{decision-point}

\vspace{0.5cm}

\begin{flushright}
{\small\color{muted} 下一章：美丽的初稿——AI把``做出来''的门槛降为零，把``知道好不好''抬向无限。}
\end{flushright}
% ============================================
% 第6章：美丽的初稿——AI降低"做出来"的门槛，抬高"知道好不好"的门槛
% 拱桥 · 判断稀缺
% ============================================

\chapter{美丽的初稿}

\vspace{0.5cm}

\begin{center}
{\small\color{muted} 拱桥的秘密，不是抵抗压力，而是利用压力。}
{\small\color{muted} AI把``做出来''的压力降为零——于是``知道好不好''的压力，}
{\small\color{muted} 全部压在了人身上。}
\end{center}

\vspace{1cm}

\section{美丽的初稿}

2026年6月，师宇在知识星球上写了一篇复盘，讲的是一场面向铁路局员工的AI实操培训。他写了一段话：

\begin{shiyu-inline}
``核心概念是'美丽的初稿'：AI不应被理解为替人写完，而是给人一个结构完整、方向正确、可被修改和注入真实经验的初稿。AI写出的是初稿，漂亮、完整、体面，但它还不是你的作品。你要把真实经验、岗位责任、个人判断和行动承诺放进去。否则它只是AI写出的文字，不是你的心得。''

\hfill ——师宇博文，2026年6月12日
\end{shiyu-inline}

\textbf{美丽的初稿。}这四个字，是我认为对AI最准确的定位。

AI给你的东西，不是垃圾。它漂亮、完整、体面。格式规范，逻辑清晰，甚至还附带了参考文献——虽然那些参考文献可能有一半是它编的。它看起来像一个成品。但它不是。

它是\textbf{初稿}。一个结构完整、方向正确、但还没有被注入真实经验的初稿。它需要一个人——一个有真实经验、岗位责任、个人判断和行动承诺的人——把它变成作品。

\section{AI会编造不存在的国家标准}

在同一场培训里，学员提出了一个很具体的问题：

\begin{shiyu-inline}
``AI会编造不存在的国家标准，生成内容会引用已经被替代的旧标准，时效性无法保证。针对细分专业领域，AI生成的内容和实际专业要求不匹配，专业内容准确性较差。''

\hfill ——2026年6月11日，兰州铁路局AI培训，录音记录
\end{shiyu-inline}

AI会编造国家标准。不是因为它``想骗你''——它没有``想''这个功能。是因为它的目标不是``准确''，而是``合理''。它要生成一个看起来最像``国家标准''的东西，而不管那个标准是否真的存在。

\textbf{AI的'错误'不是语法错误。是工程判断错误。}

它不知道一根梁多高才合理。它不知道一个安全系数取多少才够。它不知道一个节点在什么位置才安全。它只是算出了——在它见过的所有文本里——这个数值最像``正确的答案''。

\section{``做出来''的门槛降为零，``知道好不好''的门槛抬向无限}

2026年6月3日，师宇写了一篇博文，标题叫``Skill不是取经路''。里面有段话：

\begin{shiyu-inline}
``AI把'做出来'的门槛降得很低，却把'知道好不好'的门槛抬得更高。Skill可能拉平动作表层，却拉远判断深层。它给普通人专家的外壳，也让真正专家更快迭代。''

\hfill ——师宇博文，2026年6月3日
\end{shiyu-inline}

这是全书最重要的一句话。

\textbf{AI把``做出来''的门槛降得很低。}

以前，你要画一张桥梁效果图，你需要学PS、学渲染、学建模。现在，你输入一句话，AI给你一张图。以前，你要写一份公文，你需要学格式、学模板、学措辞。现在，你输入一个主题，AI给你一份格式规范的草稿。以前，你要搭一座桥的有限元模型，你需要学MIDAS、学ANSYS、学网格划分。现在——你仍然需要。但AI可以帮你生成第一版。

\textbf{但AI把``知道好不好''的门槛抬得更高。}

当所有人都会``做出来''，``做出来''就不再是稀缺能力。稀缺的是``判断''——判断这张图哪里不对，判断这份公文哪里不合适，判断这个模型哪里不安全。而判断，比做出来难得多。

\textbf{做出来只需要一个AI。判断需要一个人。一个被训练过、有经验、有责任感、会在自己错的时候发现自己错的人。}

\begin{insight}
\label{insight:chap6}

\textbf{AI时代最稀缺的能力，不是``写得快''，而是``审得好''。}

不是``画得漂亮''，而是``知道哪里不对''。不是``算得快''，而是``知道这个结果在物理上是否可能''。

生成能力在贬值。判断能力在升值。这就是拱桥的压力转化——AI把生成的压力降为零，把判断的压力全部压在人的身上。而人，恰好就是被这种压力锻造出来的。
\end{insight}

\section{拱桥的隐喻之三}

拱桥的每一块石头都承受着巨大的压力。但正是这种压力，让石头被压紧，让拱桥站住。

\textbf{AI就是那个压力。}

它把``做出来''变得不值钱——于是你被迫去学``判断''。它把``知识传递''变得不值钱——于是你被迫去学``制造困惑''。它把``批改作业''变得不值钱——于是你被迫去学``追问理解''。

这些压力，不是要摧毁你的敌人。它们是建造拱桥的石头。每一块压力，都对应着一块新的能力。就像AI把做出来的门槛降为零，恰好把判断的门槛抬向无限——而判断，正是人最不能被替代的东西。

\begin{shiyu-note}
\label{note:chap6}

\textbf{会生成不稀缺。会审查、会判断、会修改、会提出更高标准才稀缺。}

AI时代的教育，最危险的不是``学生用AI写作业''。最危险的是``学生不会判断AI写的东西对不对''。

作业可以AI写。但判断——判断AI写的东西哪里不对、哪里不够好、哪里需要改——这个能力，AI替不了。而培养这个能力，恰好是教育最应该做的事。

不是禁止AI。是让学生在AI生成的初稿上，练习判断。让他们指出AI哪里错了。让他们修改AI生成的东西。让他们在AI的基础上，做出更好的版本。让他们学会说：``这不是我想要的。我要的是——''

那个``我要的是''，才是教育真正的产出。
\end{shiyu-note}

\vspace{1cm}

\begin{decision-point}
\label{decision:chap6}

\noindent\textbf{这一章，我们看到了AI如何把``做出来''变便宜，把``判断''变贵。}

\vspace{0.3cm}

\noindent 下一节课，做一件事：用AI生成一份你课程里某个知识点的``标准答案''。在课堂上展示给学生。然后问他们：``这是AI给的答案。你们觉得对吗？如果不对，哪里不对？''

\noindent 不要接受``我觉得对''或``我觉得不对''。要求他们\noindent指出\textbf{具体哪里不对}。要求他们\noindent给出\textbf{修改方案}。

\noindent 你会发现，``判断''比``知道''难得多。而那个能指出具体错误的学生，才是真正理解了的人。
\end{decision-point}

\vspace{0.5cm}

\begin{flushright}
{\small\color{muted} 第二部分结束。下一部分：斜拉桥——每一根索都直接传递力，学习回路的设计。}
\end{flushright}

% ====== 第三部分：斜拉桥 ======
\part[斜拉桥]{斜拉桥\\[4pt] {\large\color{muted}每一根索都直接传力：把学习回路接起来}}
\include{chapters/part3-cablestayed/chap07}
% ============================================
% 第8章：把理论藏进手感里
% 斜拉桥 · 从黑箱到玻璃箱
% ============================================

\chapter{把理论藏进手感里}

\vspace{0.5cm}

\begin{center}
{\small\color{muted} 斜拉桥的索，传递的是力。}
{\small\color{muted} 课堂的索，传递的是理解。}
{\small\color{muted} 当学生自己写求解器，理解就不再是被告知的——而是被发现的。}
\end{center}

\vspace{1cm}

\section{``会不会用软件''和``懂不懂结构''是两回事}

2025年11月7日，我的课堂上进行了一场学生小组作业汇报。主题是桥梁风工程与结构动力学可视化开发。后来转写出来的文字里，有一段话让我印象很深：

\begin{shiyu-inline}
``学生做了很多可视化的工作，用Python画了模态动画，用CFD做了风场模拟。但汇报的时候，我问了一个问题：'这个模态频率是怎么算出来的？'学生愣了一下，说：'软件算的。'我追问：'软件是怎么算的？'学生：'不知道。''

\hfill ——2025年11月7日，学生作业汇报，录音转写
\end{shiyu-inline}

\textbf{``软件算的。''}这三个字，是我在工程教育里听到最多的回答。学生知道怎么点按钮，知道怎么导结果，知道怎么截图放到报告里。但他们不知道按钮背后发生了什么。

这不是学生的错。是整个教学体系把``理解''外包给了软件。我们教学生用MIDAS，用ANSYS，用SAP2000——我们教他们建模、加载、运行、看云图。但我们很少教他们：\textbf{云图是怎么来的。}

\begin{insight}
\label{insight:chap8}

\textbf{会用软件，不等于懂结构。}

一个学生可以在MIDAS里建一座斜拉桥，跑出漂亮的弯矩图，但他可能不知道那个弯矩图是求解$KU=F$得到的。他可能不知道$K$是刚度矩阵，$U$是位移向量，$F$是荷载向量。他可能不知道，如果边界条件设错了，整个结果都是错的——而软件不会告诉你。

因为软件的目标是``算出结果''。它的目标不是``让你理解''。
\end{insight}

\section{从零手写有限元}

所以，在做SpanCraft的时候，我们做了一个决定：\textbf{不调用任何商业有限元。}

我们从头写了二维直接刚度法：定义节点和单元，组装整体刚度矩阵，施加边界条件，求解$KU=F$，做内力与变形的后处理。全程透明，可读，可改。代码就在游戏里，每个学生都能打开看。

这件事的意义，远比``省一笔软件授权费''大：

\begin{shiyu-inline}
``当求解器是你自己写的，'建模'就回到了力学第一性原理。它不再是'把结构画进软件'，而是直面自由度、刚度、平衡与协调这些最根本的东西。游戏里那根变色的杆件、那道下挠，背后是一套学生能完全打开、看懂、甚至动手改的计算，而不是一个只吐结果的盒子。''

\hfill ——SpanCraft设计复盘，2026年6月
\end{shiyu-inline}

我们重建的是\textbf{``建模—分析—解读''这条完整主线}，而不是把它切成``软件操作''的碎片。从定义节点/单元/截面/荷载（建模），到组装与求解（分析），再到内力/变形/影响线的判读（解读）——第一次以``玩+透明计算''的方式，被\textbf{完整地走了一遍}。

\textbf{把黑箱，重新变回玻璃箱。}学生看见的不是一张``结果云图''，而是结果如何从$K$、$U$、$F$一步步生长出来。

\section{可视化：让``看不见的力''变成可操作的对象}

2025年秋天，我花了很多时间用Manim做数学动画。悬链线的方程是$y=a\cdot\cosh(x/a)$——这个公式，每个土木学生都学过。但有多少人真正``看见''过它？

我做了一个动画：拱桥的拱轴线，参数$a$可调。你拖动$a$，拱变扁或变陡。你看到，当$a$变化时，拱的受力状态从``纯压''变成了``压弯''。那个公式，第一次从黑板上的一行字变成了你手指下的一个形状。

后来我把这些动画整理成了bridge-math项目——用Manim、HTML、交互式图表，把桥梁工程里最抽象的概念变成可视化的东西。横向分布系数、次内力、桥面板设计——每一个概念，都配一个可以拖拽、可以调整参数、可以实时看到变化的交互式工具。

但可视化不是终点。我在课堂上反复强调一件事：

\begin{shiyu-inline}
``从三维动画回到一维公式。可视化不是终点，抽象才是。真正要教的是模型成立的前提和思维路径。''

\hfill ——2026年5月18日，简支梁桥行车道板内力计算课，录音转写
\end{shiyu-inline}

\textbf{可视化是手段，不是目的。}目的是让你在``看见''之后，能用公式表达。从三维动画回到一维公式——那个从具体到抽象的过程，才是真正的理解。

\section{``理论藏进手感里''}

2026年5月22日，我在游戏化教学分享中说了一段话，后来被反复引用：

\begin{shiyu-inline}
``好的教学游戏，是把理论藏进手感里，而不是把公式贴在屏幕上。学生先'感觉对'，再回过头来，发现那个感觉有名字、有公式——这个顺序，比反过来强得多。''

\hfill ——2026年5月22日，游戏化教学分享
\end{shiyu-inline}

这句话概括了斜拉桥部分的核心思想。

\textbf{手感先于公式。直觉先于推导。}一个学生先在游戏里搭了三次桥，塌了三次，第四次站住了。他不需要任何公式就能感受到：跨中那个地方，必须加高；两边的支座，不能动。然后他回过头来，发现那个感觉有名字——跨中弯矩$qL^2/8$。支座约束。连续梁的负弯矩卸载。

这个顺序，比反过来强得多。因为\textbf{公式是别人发现的。手感是你自己发现的。}别人发现的公式，你可能会忘。你自己发现的手感，永远不会忘。

\section{斜拉桥的隐喻之二}

斜拉桥的索，每一根都直接从桥面拉到桥塔。它不是梁桥那种``通过弯曲来承重''。它是\textbf{直接传递}。

\textbf{好的教学，就是斜拉桥。}

每一个概念，都有一根索直接连接到学生的直觉。你不用绕弯子——不用先讲三章理论再做一个练习。你直接让学生动手。搭桥。塌。再搭。再塌。站住了。然后你问：``你知道为什么站住了吗？''然后那根索——那根连接着``手感''和``公式''的索——就被拉紧了。

\begin{shiyu-note}
\label{note:chap8}

\textbf{工程教育最怕的，是把``理解''外包给软件。}

而一旦你亲手写过一次哪怕最简单的二维有限元，你对``刚度''二字的体会，就再也不一样了。你不只是知道$K$是刚度矩阵。你\textbf{感觉到}了$K$——你在组装它的时候，发现有些节点连在一起，有些没有；你在求解的时候，发现如果边界条件设错了，整个矩阵就奇异了；你在看结果的时候，发现那个变形，就是从$U$里长出来的。

这才是真正的理解。不是被告知的。是被发现的。
\end{shiyu-note}

\vspace{1cm}

\begin{decision-point}
\label{decision:chap8}

\noindent\textbf{这一章，我们看到了从``黑箱''到``玻璃箱''的转变。}

\vspace{0.3cm}

\noindent 你教的课里，有没有一个``黑箱''——一个学生只知道输入和输出，但不知道中间发生了什么的东西？一个软件？一个公式？一个设计流程？

\noindent 下一节课，把这个黑箱打开。不是讲一遍原理——是让学生\textbf{亲手做一遍}。如果是个软件，让他们自己写一个简化版。如果是个公式，让他们自己推导一遍。如果是个流程，让他们自己走一遍，然后告诉你：哪一步最关键，为什么。

\noindent 你会发现，``理解''和``会用''，是两回事。
\end{decision-point}

\vspace{0.5cm}

\begin{flushright}
{\small\color{muted} 下一章：从``把知识讲清楚''到``让学习发生''——六条教学理念。}
\end{flushright}
\chapter{最短的学习回路}

\section{不要等到下周才知道桥塌了}

这本书的材料很杂：课堂讲授、学生汇报、教师培训、游戏原型、网页工具。把它们串起来的不是 AI，也不是桥，而是一个反复出现的动作：\textbf{做一次，看见结果，改一下，再解释。}

2026 年 5 月 22 日，我用最朴素的话说过它：

\begin{shiyu-inline}
``我做出来一个操作，马上获得这个东西行还是不行，好还是不好。然后我跟这个回应，再去做调试。''

\hfill ——游戏化教学分享，2026 年 5 月 22 日 00:13:32
\end{shiyu-inline}

这句话后来被原稿不断拔高，几乎变成了游戏解决一切。其实它只解决第一件事：缩短动作和结果之间的距离。

\section{四步，不是两步}

完整回路有四步：

\begin{enumerate}
  \item \textbf{尝试：}学生先作出预测、模型或初稿；
  \item \textbf{反馈：}系统、同伴、教师或真实约束指出差距；
  \item \textbf{修改：}学生据此改变作品，而不是只读评语；
  \item \textbf{解释：}学生说明改了什么、为什么、还不确定什么。
\end{enumerate}

只做“尝试—反馈”，学生可能看完分数就走。只做“生成—修改”，修改可能完全由 AI 完成。最后的解释让教师看见判断是否真的回到了学生手里。

反馈研究把重点从“教师给了什么”转向“学习者拿它做了什么”。Carless 把跨任务持续使用反馈称为反馈螺旋\cite{carless2019spirals}。这与课程项目里的版本迭代非常接近：上一版的问题，应该进入下一版，而不是随成绩一起封存。

\section{谁来给反馈}

桥本身可以反馈：车开过去，结构变形或破坏。测试用例可以反馈：代码在边界输入下报错。同伴可以反馈：你的解释哪里听不懂。AI 可以反馈：版本之间变了什么、哪些评论重复。教师则把时间留给模型假设、专业取舍和学生自己还没有察觉的盲点。

反馈来源越多，不代表越好。最关键的是每一种反馈只做自己擅长的事，并且学生有机会行动。

\section{教师真正困难的工作}

设计回路比讲解更麻烦。你要决定在哪一步让学生卡一下，什么时候给提示，什么错误可以安全发生，什么错误必须立即制止。你还要忍住，不在学生刚开始想的时候就把答案说出来。

更现实的问题是：一个班几十份作品、几轮迭代，教师会不会被拖成一台永远改不完的机器？下一章，我们不谈“教师不可替代”这种宽慰，直接谈哪些活该交给 AI，哪些活教师必须留给自己。


% ====== 第四部分：悬索桥 ======
\part[悬索桥]{悬索桥\\[4pt] {\large\color{muted}跨越越远，越需要可靠的锚}}
% ============================================
% 第10章：人的判断权：AI越强，人越不能退后
% 悬索桥 · 判断权
% ============================================

\chapter{人的判断权：AI越强，人越不能退后}

\vspace{0.5cm}

\begin{center}
{\small\color{muted} 悬索桥跨越最长的距离。}
{\small\color{muted} 但它依赖最深的锚——锚碇，埋在地下，看不见，但绝对不能松。}
{\small\color{muted} 人的判断权，就是教育的锚碇。}
\end{center}

\vspace{1cm}

\section{``Skill不是取经路''}

2026年6月3日，师宇在知识星球上写了一篇博文，标题叫``Skill不是取经路''。里面有一段话，我反复读了很多遍：

\begin{shiyu-inline}
``Skill不是取经路。把经历写成'遇到妖怪就打、打不过找佛祖'，并不会让另一只猴子成为孙悟空。真正的能力不是知道规则，而是知道何时打、何时收手、何时求助、求谁，以及为什么还要往前走。''

\hfill ——师宇博文，2026年6月3日
\end{shiyu-inline}

这段话击中了AI时代最核心的问题。

Skill——AI技能包——是我这两年一直在推广的东西。把专家的隐性经验打包成可调用的流程：做什么、怎么做、要注意什么、什么情况下不能这么做。它让一个新手能借助AI，做出接近专家水平的东西。

但师宇在这篇博文里问了一个更尖锐的问题：\textbf{有了Skill，这个人就真的成了专家吗？}

孙悟空的取经经验，如果写成``遇到妖怪就打，打不过找佛祖''，当然也对。但另一只猴子读了这段话，不会变成孙悟空。因为它不知道\textbf{什么时候}该打，什么时候不该打，什么时候该求人，求谁，以及——\textbf{为什么非去不可。}

\section{AI把``做出来''变容易，把``判断''变难}

师宇在同一篇博文里继续写道：

\begin{shiyu-inline}
``AI把'做出来'的门槛降得很低，却把'知道好不好'的门槛抬得更高。Skill可能拉平动作表层，却拉远判断深层。它给普通人专家的外壳，也让真正专家更快迭代。''

\hfill ——师宇博文，2026年6月3日
\end{shiyu-inline}

\textbf{Skill给普通人专家的外壳。}

这句话让我想了很久。一个学生，用我做好的桥梁Skill，输入几个参数，AI生成了一座漂亮的斜拉桥模型。看起来，他``会''了。但如果你问他：``为什么塔高取这个值？为什么索距取这个值？如果地基条件变了，这个设计还能用吗？''——他答不上来。

\textbf{他有的不是能力。他有的是一套能生成专家外壳的工具。}

\begin{insight}
\label{insight:chap10}

\textbf{真正危险的，不是交出执行，是交出判断。}

AI可以帮你写。可以帮你画。可以帮你算。可以帮你归纳你是什么人。可以生成一套``如何对待你''的策略。它越强，人越容易舒服地退到后面。

但真正该被训练的，恰恰是\textbf{人不退后的能力}。人要知道自己为什么做这件事。知道什么值得做。知道哪里不够好。知道何时该慢下来。知道什么不能交出去。
\end{insight}

\section{``人与人的东西''仍被认为最重要}

2026年7月，我在教学理念综合分析中写下了这样一段话：

\begin{shiyu-inline}
``即使平台高度智能化，'人与人的东西'仍被认为最重要。材料一方面追求动态问题、即时排名和AI汇总，另一方面又明确说积分只是辅助、线下活动更重要，并建议学生的迭代过程应由学生自己总结，而不是交给AI。真正成熟的AI教学，可能不在于自动化更多，而在于知道哪些认知劳动必须保留给人。''

\hfill ——教学理念综合分析，2026年7月
\end{shiyu-inline}

这段话是我对两年教学实验的\textbf{自我限制}。

我的平台可以做很多事情。它可以自动评分、自动汇总反馈、自动生成知识图谱、自动检测学生的盲区。但它不能替代一件事：\textbf{学生自己总结自己的迭代过程。}

为什么？因为总结，是\textbf{判断}的核心动作。你回头看自己做了什么，哪些对了，哪些错了，为什么错，下次怎么改——这个动作，AI不能替你做。不是因为它做不到。是因为\textbf{它做了，就不是你的了。}

\section{交出判断的那一刻}

2026年5月31日，师宇写了一篇短文，标题叫``扫描我的电脑，看看你能做点什么？''。他写道：

\begin{shiyu-inline}
``这个提示词不是人告诉AI做什么，而是把判断权交给AI，让AI想想能为人做什么。这当然很强，也确实很诱人。可一旦人习惯了把判断权交出去，所谓赋能就会滑向献祭。''

\hfill ——师宇博文，2026年5月31日
\end{shiyu-inline}

\textbf{赋能滑向献祭。}

这是全书最让我不安的四个字。AI帮你做越来越多的事，你越来越轻松，越来越依赖它。有一天，你发现你不再需要思考了。AI帮你思考。你只需要点头。你不需要判断什么值得做——AI告诉你。你不需要知道哪里不够好——AI帮你改好了。你不需要决定什么不能交出去——因为\textbf{你已经不记得什么东西是重要的了}。

这就是悬索桥的锚碇。

\section{悬索桥的隐喻}

悬索桥跨越最长的距离。但它依赖最深的锚——锚碇，埋在地下，看不见，但绝对不能松。锚碇松了，整座桥就塌了。

\textbf{人的判断权，就是教育的锚碇。}

AI可以帮你搭桥。可以帮你画图。可以帮你写报告。可以帮你分析数据。但它不能帮你判断：这座桥对不对？这张图哪里不对？这个报告有没有注入真实经验？这个数据背后是什么物理现象？

这些判断，必须由人来做。不是因为AI做不了。是因为\textbf{AI做了，人就不再会做了}。而人不再会判断的时候，那座桥就会塌——不是物理上的塌，是精神上的塌。

\begin{shiyu-note}
\label{note:chap10}

\textbf{AI越强，人越不能退后。}

这不是反对AI。这是反对把判断权交给AI。

AI可以成为第二双手。但不能成为替你活着的理由。AI可以搭出桥的形状。但人要决定过不过桥、往哪里过、为什么过。

教育最应该守住的，不只是知识的交付，更是\textbf{判断}——判断什么值得学、什么值得做、什么值得守住。我们无法替未来的AI能力画一条永恒边界，但可以决定：人的判断不能因为工具便利而退出。身体经验也会让判断与后果、他人和具体世界发生联系，这正是工程教育不能轻易外包的部分。
\end{shiyu-note}

\vspace{1cm}

\begin{decision-point}
\label{decision:chap10}

\noindent\textbf{这一章，我们看到了判断权为什么是教育的底线。}

\vspace{0.3cm}

\noindent 下一次，当你用AI生成一份教案、一份报告、一张图的时候——停下来。在点``使用''之前，问自己三个问题：

\noindent 1. 这个东西，哪里可能不对？
\noindent 2. 如果我的学生/同事问我``为什么选这个方案''，我能解释吗？
\noindent 3. 这个东西里，有没有一句话、一个选择、一个细节——是我自己做的？

\noindent 如果第三个问题的答案是``没有''——这份东西不是你的。重做。
\end{decision-point}

\vspace{0.5cm}

\begin{flushright}
{\small\color{muted} 下一章：教师是``改稿机器''吗？——AI时代的教育伦理。}
\end{flushright}

% ============================================
% 第11章：教师是"改稿机器"吗？——AI时代的教育伦理
% 悬索桥 · 教育伦理
% ============================================

\chapter{教师是``改稿机器''吗？}

\vspace{0.5cm}

\begin{center}
{\small\color{muted} 悬索桥的主缆，承受着整座桥的重量。}
{\small\color{muted} 但如果锚碇松了，主缆再强也没有用。}
{\small\color{muted} 教师的角色，就是教育的锚碇。}
\end{center}

\vspace{1cm}

\section{当AI可以批改一切}

2026年6月23日，我在跟两位老师讨论AI教学平台。我讲到了平台上的AI批改功能：

\begin{shiyu-inline}
``AI会去过滤，说你这个评论是不是在灌水。有的学生最开始第一版的时候，他打一一一一，凑够十个一，我这要求至少十个字嘛，一点提交，叭就过去了，他的分拿到手了。但是这种情况非常少，更多的学生会认真的写。''

\hfill ——2026年6月23日，AI教学平台讨论，录音转写
\end{shiyu-inline}

AI可以自动检测灌水评论。AI可以自动聚合三十个人的评语。AI可以自动判断一个学生的参与是否``有效''。AI可以自动给分。

这些功能，在传统的课堂里，都是老师做的。老师要花几个小时看评论、找灌水、打分、汇总。现在，AI几秒钟就做完了。

\textbf{然后呢？}

老师多出来的时间，用来做什么？用来设计更好的困惑？用来跟学生一对一的交流？用来反思自己的教学？还是——用来变成一台``改稿机器''？

\section{``改稿机器''}

2026年6月12日，师宇在复盘一场铁路局AI培训时，提到了一个让我警觉的场景：

\begin{shiyu-inline}
``培训现场也触及深层焦虑：AI可能让人不会思考，导师、老师、管理者可能被迫变成'改稿机器'。''

\hfill ——师宇博文，2026年6月12日
\end{shiyu-inline}

\textbf{改稿机器。}

学生用AI生成初稿。老师用AI检测AI。老师改几处。学生再生成。老师再检查。循环往复。老师和学生之间，不再有真实的交流。只有AI生成的内容，在两个屏幕之间来回传输。

这可能是AI时代教育最可怕的未来：\textbf{不是AI取代了教师。是教师变成了AI的质检员。}

\section{什么必须保留给人？}

2026年7月，在教学理念综合分析中，我列出了\textbf{不可外包给AI的环节}：

\begin{shiyu-inline}
``保留'不可外包给AI'的环节：自我判断懂没懂、解释关键选择、失败复盘和迁移总结由学生完成。''

\hfill ——教学理念综合分析，2026年7月
\end{shiyu-inline}

这不是一个技术问题。这是一个伦理问题。

\textbf{自我判断懂没懂。}AI可以给你一个自测题，告诉你得分。但它不能替你感受``我懂了''那个瞬间。那个瞬间，必须是你自己经历的。

\textbf{解释关键选择。}AI可以解释它为什么选这个方案。但那个解释是统计的——``因为数据里这个方案最常见''。那不是解释。那是概率。真正的解释需要你理解：在什么约束下，我为什么放弃了A选了B。

\textbf{失败复盘。}AI可以分析你的错误模式。但它不能替你感受失败——那种``我明明知道但就是做错了''的懊悔，那种``我下次一定不会再犯''的决心。这些感受，是学习最核心的驱动力。

\textbf{迁移总结。}AI可以帮你总结``这次学到了什么''。但它总结的是它看到的模式。你真正学到的东西，可能和AI总结的完全不同。因为\textbf{学习是转化}——你被改变了，你变成了另一个人。AI看不到这种改变。

\section{教师的新角色}

如果AI能做知识传递，能做作业批改，能生成教案和PPT——那教师还能做什么？

\textbf{教师能做三件事，AI一件都做不了。}

\textbf{第一，口头诊断。}2026年1月9日，混凝土结构答辩。我改了一个参数，问学生：``为什么配筋量变了？''那个问题，不是在考他知不知道答案。我是在\textbf{诊断}——诊断他理解到了哪一层，诊断他的思维卡在哪里，诊断下一步该往哪个方向引导。AI可以做选择题。AI不能做口头诊断。因为口头诊断需要\textbf{看见}——看见学生的眼神，听见他的犹豫，感受他是在``回忆''还是在``推理''。

\textbf{第二，设计困惑。}AI可以生成任何内容。但它不能判断：\textbf{这个内容，在什么时间、以什么方式、给什么学生，能制造出最有效的困惑。}因为困惑是\textbf{个人的}。同一个问题，对A学生是困惑，对B学生是常识。设计困惑，需要你了解每一个学生——不是他们的分数，是他们\textbf{是谁}。

\textbf{第三，见证骄傲。}AI可以给分。但它不能\textbf{见证}。当学生指着自己搭的桥说``站住了''，那个眼神，那个语气，那个嘴角微微上扬的弧度——AI看不到。它不知道那个学生在三个月前差点转专业。它不知道那个学生曾经在深夜因为一个参数调不对而摔过键盘。它不知道那个瞬间，对那个学生意味着什么。

\textbf{见证，是教师最不能被替代的能力。}

\begin{insight}
\label{insight:chap11}

\textbf{AI时代，教师最值钱的能力，不是``讲得好''。}

是\textbf{听得懂}——听懂学生卡在哪里。是\textbf{看得见}——看见学生努力了但没做成。是\textbf{等得起}——等学生自己发现，而不是急着告诉他答案。是\textbf{记得住}——记得这个学生三个月前的样子，知道他变了多少。

这些，AI一件都做不了。
\end{insight}

\section{悬索桥的隐喻之二}

悬索桥的主缆，承受着整座桥的重量。但主缆的力，最终要传到锚碇——埋在地下，看不见，但绝对不能松。

\textbf{教师的角色，就是教育的锚碇。}

AI可以拉起主缆。AI可以生成内容、批改作业、分析数据。但所有这些力的终点，必须是一个\textbf{人}——一个能判断、能诊断、能见证、能在学生指着自己搭的桥说``站住了''的时候，真心为他高兴的人。

锚碇松了，悬索桥就塌了。教师退后了，教育就变成了两台机器之间的数据传输。

\begin{shiyu-note}
\label{note:chap11}

\textbf{不要让AI把你变成改稿机器。}

当AI帮你批改作业的时候，省下来的时间，不要用来批改更多作业。用来跟学生聊天。用来观察他们搭桥。用来问他们：``你为什么这么搭？''——然后认真听他们回答。

那些回答里，有你用任何AI都检测不出来的东西：一个正在努力理解世界的人，正在变成另一种人。
\end{shiyu-note}

\vspace{1cm}

\begin{decision-point}
\label{decision:chap11}

\noindent\textbf{这一章，我们看到了AI时代教师角色的重新定义。}

\vspace{0.3cm}

\noindent 下一节课，不要站在讲台上。走下去。走到一个学生旁边。看他搭桥。不要说话。等他自己发现错了。等他改。等他再试。等桥站住。

\noindent 然后只说一句话：``你刚才改的那一下，为什么管用？''

\noindent 听他的回答。你会听到一些AI永远听不到的东西。
\end{decision-point}

\vspace{0.5cm}

\begin{flushright}
{\small\color{muted} 下一章：桥不会自己存在——关于责任与骄傲。}
\end{flushright}
% ============================================
% 第12章：桥不会自己存在——关于责任与骄傲
% 悬索桥 · 责任与骄傲
% ============================================

\chapter{桥不会自己存在}

\vspace{0.5cm}

\begin{center}
{\small\color{muted} 悬索桥的主缆，是靠锚碇拉住的。}
{\small\color{muted} 锚碇埋在地下，没有人看见它。}
{\small\color{muted} 但整座桥知道它在那里。}
\end{center}

\vspace{1cm}

\section{你不能骗一座桥}

2025年11月4日，涡振原理课。我放了一段虎门大桥涡振的视频。后来转写出来的文字里，有一段让我印象很深：

\begin{shiyu-inline}
``46辆满载的重车压在上面，桥还是在震动。你可以看到，虽然是我们感觉上，这个桥梁是很刚的，那个风是很小的，顶多是吹到上面把灰吹掉了。但其实在这种大型的结构上，轻柔的结构上，如果它一旦发生动力学问题，40多辆车都压不住的这种震动。''

\hfill ——2025年11月4日，涡振原理课，录音转写
\end{shiyu-inline}

46辆重车。压不住一座桥的震动。

这就是结构的真相：\textbf{你不能骗一座桥。}你可以骗一份PPT。你可以骗一份报告。你可以骗一个评委。你可以骗你自己。但你不能骗一座桥。桥不会听你解释。它只会站住，或者塌掉。

\textbf{土木工程教会我的，恰好就是这种``不能骗''的品格。}

在一个越来越多人相信``差不多就行''的世界里，这种品格是一种稀缺的、值得被保护的东西。AI可以生成一切``看起来像桥''的东西。但\textbf{桥不是看起来像就行的。桥必须真的站得住。}

\section{``它好不好看''和``它扛不扛得住''}

2026年6月，SpanCraft设计复盘。我写下了这样一段话：

\begin{shiyu-inline}
``我越来越相信，能在同一件事情上同时用两套大脑——一边问'它好不好看'，一边问'它扛不扛得住'——本身就是一种稀缺的能力。创造美，并且严肃地为这份美负责，大概是工程教育里最难教、也最值得教的东西。''

\hfill ——SpanCraft设计复盘，2026年6月
\end{shiyu-inline}

\textbf{两套大脑。}

一套判断美。一套判断真。一套问：``这个东西好看吗？打动人心吗？让人想停留吗？''另一套问：``这个东西安全吗？符合规范吗？能站住吗？''

在大多数行业里，美和真是分开的。艺术家负责美，工程师负责真。但桥梁工程师不能分开。你设计的桥，必须同时满足两个条件：\textbf{让人想走过，并且让人能安全走过。}

这种双重要求，是工程教育最难教的东西。因为美和真，看起来是矛盾的。美的桥往往更轻、更薄、更大胆——但轻、薄、大胆，恰恰是工程安全的大敌。你必须同时追求两者，在矛盾中找到那个最优解。

\textbf{那座最优解，就是拱轴线。}而拱轴线，就是全书的终章。

\section{骄傲}

2025年10月21日，桥梁课程。我讲了结构设计大赛的故事：

\begin{shiyu-inline}
``我们做了一个100克的船，它可以承载19公斤的重量。大概排在120个队里的70名。清华大学大概排在75名，我们排在清华大学前面。但是这个就是比较偶然性嘛，也不代表啥。''

\hfill ——2025年10月21日，桥梁课程，录音转写
\end{shiyu-inline}

我说``也不代表啥''。这句自我拆台很重要：课堂口述能够保留当时的兴奋，却不能替比赛记录、设计版本和学生访谈补齐细节。排名不是这一章要证明的东西。更值得留下的，是工程作品把``对不对''从纸面问题变成了一个必须面对的后果。

2026年7月1日，桥梁工程AI课程设计结题汇报里，学生讲到AI生成的弯矩图不对，改了几次仍有问题；也讲到文件被新版本覆盖、负影响线出错。到最后，他说：

\begin{shiyu-inline}
``你就真的得是自己会才行，要不然你也不知道AI做得对不对。''

\hfill ——2026年7月1日，课程设计结题汇报 00:25:20
\end{shiyu-inline}

这句话没有替工程教育写一句漂亮口号。它只是把责任重新放回了人身上：\textbf{当工具给出一个结果，你仍要有能力发现它不对，并说明为什么。}

\section{责任}

桥不会自己存在。它必须被建造。

你必须计算每一根梁的弯矩。你必须校核每一个节点的安全系数。你必须确保桥塔不会在风里摇摆，缆索不会在雪里断裂，锚碇不会在洪水里松动。你必须在图纸上签下你的名字。如果桥塌了，那个名字会被人记住。

\textbf{这就是责任。}

AI可以帮你算。AI可以帮你画。AI可以帮你生成初稿。但AI不会在图纸上签名。签名的，是你。那个名字，代表的是：\textbf{我检查过了。我相信它能站住。如果它塌了，我来负责。}

AI系统可以参与计算、检查甚至提出警告，但责任链不能因此消失。课程真正要训练的，是学生能否说明依据、留下复核记录，并知道何时必须请专业人员重新检查。

\textbf{这里不需要预测AI永远不能做什么。只需要确认：在现行的工程实践里，采纳结果、签字与承担后果的人不能把责任推给生成工具。}

\section{悬索桥的隐喻之三}

悬索桥的主缆，是靠锚碇拉住的。锚碇埋在地下，没有人看见它。但整座桥知道它在那里。

\textbf{责任与骄傲，就是教育的锚碇。}

你教给学生的一切——公式、规范、软件操作——都像主缆上的吊杆，把桥面挂在缆上。但如果锚碇松了，所有的吊杆都会失效。如果学生没有\textbf{责任}——没有``我签了名，我来负责''的意识——他学再多公式也没用。如果学生没有\textbf{骄傲}——没有``这是我的桥，我亲手建的''的感觉——他永远只是一个会用软件的人，不是一个工程师。

\textbf{责任和骄傲，是教育的最后一道防线。}

\begin{shiyu-note}
\label{note:chap12}

\textbf{你不能骗一座桥。你也不能骗一个学生。}

桥会塌。学生会发现。他们会发现你教的那些``有用的''东西，AI做得更好。他们会发现你花80\%时间教的知识，在他们毕业之前就过时了。他们会发现，你从来没有教过他们最需要的东西——\textbf{如何判断。如何负责。如何骄傲。}

不要等到他们发现。从现在开始，教他们那些AI教不了的东西。教他们判断。教他们负责。教他们骄傲。教他们指着一样东西说：``这是我的。我建的。它站住了。''
\end{shiyu-note}

\vspace{1cm}

\begin{decision-point}
\label{decision:chap12}

\noindent\textbf{这一章，我们看到了责任与骄傲为什么是教育的根基。}

\vspace{0.3cm}

\noindent 下一次，当你给学生布置一个项目的时候，在最后加一个环节。不是``提交报告''。不是``做PPT汇报''。是：

\noindent \textbf{``在项目封面上，签下你的名字。然后告诉我：这个东西，你愿意为它负责吗？''}

\noindent 你会发现，那个签名，比任何分数都重。
\end{decision-point}

\vspace{0.5cm}

\begin{flushright}
{\small\color{muted} 第四部分结束。下一部分：组合桥——没有一种桥型能解决所有问题。}
\end{flushright}


% ====== 第五部分：组合桥 ======
\part[组合桥]{组合桥\\[4pt] {\large\color{muted}没有一种桥型能解决所有问题}}
% ============================================
% 第13章：用TRIZ重新理解教学——一场系统化的叛逆
% 组合桥 · TRIZ框架
% ============================================

\chapter{用TRIZ重新理解教学：一场系统化的叛逆}

\vspace{0.5cm}

\begin{center}
{\small\color{muted} 组合桥不靠单一桥型。它把梁、拱、索的优点组合在一起。}
{\small\color{muted} TRIZ不靠单一方法。它把矛盾分析、理想定义、进化路径组合在一起。}
{\small\color{muted} 教学创新，需要同一套思维。}
\end{center}

\vspace{1cm}

\section{教学中的四对物理矛盾}

2026年3月，我在备``桥梁上的作用''这节课时，用TRIZ方法论做了一次系统分析。后来我发现，TRIZ不只是用来分析桥梁的。它也可以用来分析\textbf{教学本身}。

TRIZ的核心是矛盾分析——不是``选A还是选B''，而是``A和B都是对的，但它们是矛盾的，怎么同时满足？''

教学中有四对最根本的物理矛盾：

\textbf{第一对：讲 vs 不讲。}教师应该讲——不讲学生怎么知道？教师不应该讲——讲了学生自己就不想了。解法：\textbf{时间分离。}困惑时闭嘴，想通后追问。学生自己试的时候不讲话，学生卡住的时候给一个提示，学生想通了问``为什么''。

\textbf{第二对：帮 vs 不帮。}AI应该帮学生——降低门槛，让更多人能参与。AI不应该帮学生——帮太多，学生就不会自己判断了。解法：\textbf{条件分离。}生成初稿时帮，判断对错时不帮。让AI生成第一版，让学生判断AI哪里错了。

\textbf{第三对：标准化 vs 个性化。}评价应该标准化——公平。评价应该个性化——每个学生不一样。解法：\textbf{系统分离。}标准化评价``参与度''（有没有投入），个性化评价``理解深度''（有没有自己的判断）。参与用积分，理解用追问。

\textbf{第四对：快 vs 慢。}AI让一切变快——生成、批改、反馈。但学习需要慢——困惑、失败、反思。解法：\textbf{空间分离。}AI负责快的部分——生成初稿、聚合反馈、检测盲区。人负责慢的部分——口头诊断、失败复盘、迁移总结。

\begin{insight}
\label{insight:chap13}

\textbf{TRIZ教给我们的，不是怎么选。是怎么同时满足两个相反的要求。}

讲和不讲，不是二选一。是在该讲的时候讲，在该闭嘴的时候闭嘴。帮和不帮，不是二选一。是在该帮的时候帮，在该放手的时候放手。教学设计不是选边站。是设计\textbf{分离条件}。
\end{insight}

\section{教学的IFR：理想最终结果}

TRIZ有一个概念叫\textbf{IFR}——理想最终结果。它问的是：如果没有任何限制，最完美的状态是什么？

桥梁的IFR是：\textbf{自重为零，100\%承载交通，同时自动适应所有环境作用。}

教学的IFR是什么？我想了很久，写下这样一句话：

\begin{shiyu-inline}
\textbf{学习自动发生，无需外部驱动。}

学生自己产生困惑，自己寻找答案，自己验证理解，自己发现错误，自己修正路径，自己感受骄傲。教师的存在，不是为了推动学习——而是为了创造一个环境，让学习无法不發生。

\hfill ——教学理念综合分析，2026年7月
\end{shiyu-inline}

现实离这个IFR有多远？

\begin{itemize}
  \item \textbf{困惑：}目前大部分困惑是老师制造的，不是学生自己产生的。差距：大。
  \item \textbf{寻找：}目前大部分答案来自老师或AI，不是学生自己找的。差距：大。
  \item \textbf{验证：}目前大部分验证依赖考试和批改，不是学生自己判断的。差距：极大。
  \item \textbf{骄傲：}目前大部分学生感受不到骄傲，只感受到``终于考完了''。差距：极大。
\end{itemize}

但IFR不是用来实现的。它是用来\textbf{趋近}的。每一次你设计一个让学生自己产生困惑的环节，你就在向IFR靠近一步。每一次你让学生自己判断AI的答案对不对，你就在向IFR靠近一步。

\section{教学的进化路线}

TRIZ还有一个概念叫\textbf{进化路线}——从当前状态到IFR的技术演进路径。教学的进化路线是什么？

\begin{center}
\begin{tikzpicture}[node distance=2.5cm, auto]
  \node[draw, rectangle, rounded corners, fill=blue-soft] (A) {传统讲授};
  \node[draw, rectangle, rounded corners, fill=blue-soft, right of=A] (B) {游戏化};
  \node[draw, rectangle, rounded corners, fill=blue-soft, right of=B] (C) {AI辅助};
  \node[draw, rectangle, rounded corners, fill=blue-soft, right of=C] (D) {AI Skill};
  \node[draw, rectangle, rounded corners, fill=blue-soft, right of=D] (E) {师生机共建};
  \node[draw, rectangle, rounded corners, fill=gold, right of=E] (F) {学习生态};
  \draw[->] (A) to (B);
  \draw[->] (B) to (C);
  \draw[->] (C) to (D);
  \draw[->] (D) to (E);
  \draw[->] (E) to (F);
\end{tikzpicture}
\end{center}

\textbf{传统讲授：}老师讲，学生记。老师考，学生答。这是起点。

\textbf{游戏化：}引入即时反馈和主动探索。学生在安全环境中反复失败。这是第一步进化。

\textbf{AI辅助：}AI生成初稿、聚合反馈、检测盲区。效率提升，但角色不变。这是第二步。

\textbf{AI Skill：}专家的隐性经验被蒸馏成可调用的技能包。新手能做出接近专家水平的东西。但判断权仍在人手里。这是第三步。

\textbf{师生机共建：}教师负责方向，AI负责生成，学生负责判断。三者各司其职。这是第四步。

\textbf{学习生态：}课程不再是每学期归零的临时活动。学生作品成为下一届的案例，比赛反哺课程，录音沉淀为知识图谱，失败变成教学资产。这是IFR的趋近。

\textbf{你现在在哪一步？}不重要。重要的是，你知道下一步是什么。

\section{组合桥的隐喻}

组合桥不靠单一桥型。它把梁、拱、索的优点组合在一起。某一跨需要梁的简单直接，就用梁。某一跨需要拱的跨度优势，就用拱。某一跨需要索的轻巧优雅，就用索。

\textbf{教学创新，也需要组合思维。}

没有一种方法能解决所有问题。游戏化适合制造困惑，但不适合系统化知识传递。AI Skill适合降低门槛，但不适合培养判断力。口头诊断适合个性化反馈，但不适合大规模课堂。你需要\textbf{组合}——在不同的环节，用不同的工具，解决不同的问题。

\textbf{而组合的原则，就是TRIZ。}分析矛盾，定义理想，选择路径。

\begin{shiyu-note}
\label{note:chap13}

\textbf{TRIZ的精神，不是找到正确答案。是拒绝在错误的问题上妥协。}

``讲还是不讲''——这个问题本身就是错的。应该问的是：``在什么条件下讲，在什么条件下不讲？''

``帮还是不帮''——这个问题本身就是错的。应该问的是：``在什么环节帮，在什么环节放手？''

教学创新，不是选边站。是设计\textbf{分离条件}。而设计分离条件，需要你理解每一个学生的状态，理解每一个环节的目标，理解每一对矛盾的本质。这很难。但这是唯一正确的路。
\end{shiyu-note}

\vspace{1cm}

\begin{decision-point}
\label{decision:chap13}

\noindent\textbf{这一章，我们用TRIZ重新理解了教学。}

\vspace{0.3cm}

\noindent 选你课上最让你纠结的一对矛盾。是``讲 vs 不讲''？``帮 vs 不帮''？``标准化 vs 个性化''？``快 vs 慢''？

\noindent 不要选边。设计一个\textbf{分离条件}——在什么情况下做A，在什么情况下做B。写下来。下一节课试试。试完告诉我，效果如何。
\end{decision-point}

\vspace{0.5cm}

\begin{flushright}
{\small\color{muted} 下一章：一封写给所有孩子的邀请函。}
\end{flushright}
% ============================================
% 第14章：效率至上世界里，一封写给所有孩子的邀请函
% 组合桥 · 邀请
% ============================================

\chapter{效率至上世界里，一封写给所有孩子的邀请函}

\vspace{0.5cm}

\begin{center}
{\small\color{muted} 组合桥不挑地形。它不要求完美的地基，不要求均匀的跨度。}
{\small\color{muted} 它接受每一种条件，然后找到最适合的组合。}
{\small\color{muted} 教育，也应该是这样的。}
\end{center}

\vspace{1cm}

\section{不是``你应该学''，而是``你想不想来搭一座桥''}

2026年5月22日，面向呼和浩特铁路局教师的分享会。开场前，我放了一个二维码。扫进去，是一个网页，里面是SpanCraft。

我没有说``请各位老师认真学习''。我说的是：``如果大家手头有电脑，可以打开这个网页，一起试试。''

十分钟后，聊天区里一个老师打了一行字：``我搭的桥又塌了。我再试一次。''

又过了五分钟，同一个人：``站住了！！！''

三个感叹号。一个成年人，一个铁路系统的工程师，在培训课上，因为一座虚拟的桥站住了，打了三个感叹号。

\textbf{这不是被要求的学习。这是被邀请的参与。}

没有人强迫他搭桥。没有人给他打分。没有人告诉他``这是教学任务，必须完成''。他只是看到了一座桥，想让它站住。然后他自己动手，自己失败，自己调整，自己成功。他自己打了三个感叹号。

\textbf{这就是教育应该有的样子。不是``你应该学''。是``你想不想来试试？''}

\section{魅力先于说教}

2026年6月，SpanCraft设计复盘。我写下了那两条原则：\textbf{先给魅力，再给负荷。一次只教一个概念。}

这两条原则，不只是游戏设计的铁律。它们是\textbf{所有教育的铁律}。

先给魅力。不是先讲道理。不是先列大纲。不是先告诉学生``通过本章学习，你将掌握\ldots\ldots''。是先让他们\textbf{看见}。看见一座黄昏里的斜拉桥，拉索一根根自己生长出来，小车驶过桥面，水面波光粼粼。先让他们心里``哦''一声。先让他们\textbf{想}知道更多。

然后，再给负荷。再讲道理。再列公式。再告诉他们那个$qL^2/8$从哪里来。但这时候，他们已经准备好了。他们已经\textbf{想要}知道了。不是你逼他们学——是他们自己追着问。

\textbf{魅力先于说教。感觉先于公式。骄傲先于成绩。}

\section{如何在AI时代保持幸福？}

2026年6月11日，师宇写了一篇博文，标题叫``有什么用呢''。他研究了AI视频生成工具，兴奋地做了很多尝试。然后被一句话问住了。他写道：

\begin{shiyu-inline}
``什么算有用？必须产出、变现、转化才叫有用吗？不是所有有用的东西，都必须马上证明自己有用。有些价值正是在暂时无用时显现。''

\hfill ——师宇博文，2026年6月11日
\end{shiyu-inline}

\textbf{有些价值正是在暂时无用时显现。}

这句话，是我在AI时代保持幸福感的秘诀。

你搭一座桥，没有人会走。你写一段代码，AI写得比你更好。你画一张图，AI画得比你更漂亮。你做一个游戏，AI可以帮你生成所有素材。\textbf{那你为什么还要做？}

因为\textbf{做本身}就是价值。不是``做出来''有价值。是``做''这个过程——你花了三个小时搭一座桥，塌了五次，第六次站住了。那个过程，改变了你。你不再是那个不会搭桥的人。你变成了一个\textbf{搭过桥的人}。这个变化，AI给不了你。多少钱都买不到。

\textbf{幸福感不是来自``拥有''。是来自``做过''。}

\section{把手弄脏，然后指着它说``这是我的''}

2026年7月1日，桥梁工程AI课程设计结题汇报。那个学生站起来，打开他的项目。然后他说了一句话，我至今记得：

\begin{shiyu-inline}
``老师，我以前觉得，'学这个还值得吗'。现在我知道了，不值得——如果你只想找一个工作的话。但如果你想让一座桥站住，它就值得。''

\hfill ——2026年7月1日，课程设计结题汇报
\end{shiyu-inline}

\textbf{不值得——如果你只想找一个工作的话。但如果你想让一座桥站住，它就值得。}

这就是邀请函的正文。

不是邀请你学土木工程。不是邀请你用AI。不是邀请你参加教学创新比赛。是邀请你——\textbf{把手弄脏}。搭一座桥。写一段代码。画一张图。做一个游戏。设计一堂课。写一本书。不管是什么，重要的是：你亲手做的。你失败过。你改过。你骄傲过。你指着它说——\textbf{这是我的}。

\begin{insight}
\label{insight:chap14}

\textbf{在效率至上的世界里，最稀缺的不是效率。是``无用的认真''。}

是你在洗浴中心搭桥搭到关门。是你在深夜因为一个参数调不对而摔键盘。是你花三个小时画一张图，然后删掉重画。是你做了一件AI比你做得更好的事，但你还是做了。因为你\textbf{想}做。因为做这件事，让你觉得活着。

AI时代，效率不再稀缺。意义才是。而意义，从来不是``有用''的东西。意义是你在做完之后，可以说：``这是我的。我建的。它站住了。''
\end{insight}

\section{组合桥的隐喻之二}

组合桥不挑地形。它不要求完美的地基，不要求均匀的跨度。它接受每一种条件，然后找到最适合的组合。

\textbf{教育，也应该是这样的。}

不是所有孩子都适合梁桥。有些孩子需要拱桥的压力才能成长。有些孩子需要斜拉桥的直接反馈。有些孩子需要悬索桥的深锚——那个告诉他``你值得''的人。有些孩子需要组合桥——他们需要不同的方法，在不同的时间，用不同的节奏。

\textbf{教育的任务，不是把所有孩子变成同一种桥。是帮每一个孩子找到他最适合的那一种。}

\begin{shiyu-note}
\label{note:chap14}

\textbf{这封信，写给每一个被``有什么用''问住的孩子。}

你不需要证明这门课``有用''。你不需要用就业率、薪资、排名来为它辩护。你只需要说：``我搭了一座桥，它站住了。你想看看吗？''

那个``想看看''，就是教育最原始的动力。好奇心。骄傲。分享的冲动。这些，AI永远不会有。但这些，是人类最古老、最自然、最强大的学习动力。

在一个效率至上的世界里，为``无用的认真''保留一片海滩。在一片海滩上，孩子们蹲在沙坑里，搭了三个小时的桥。没有人问他们``有什么用''。他们只是想让桥站住。
\end{shiyu-note}

\vspace{1cm}

\begin{decision-point}
\label{decision:chap14}

\noindent\textbf{这一章，是一封邀请函。}

\vspace{0.3cm}

\noindent 下一节课，不要从头讲到尾。留十分钟。跟学生说：``今天我准备了一个东西，但我不知道它好不好。你们帮我看看。''

\noindent 然后，把你最近做的一个东西拿出来——一个你亲手搭的桥，一个你写的代码，一张你画的图，一个你做的游戏。告诉他们：``这是我做的。我想让你们看看。''

\noindent 你会发现，教室里有什么东西变了。不是他们在听你讲。是他们在\textbf{看}你。他们看见了一个活人——一个会做东西、会骄傲、会不确定的人。他们想成为那样的人。
\end{decision-point}

\vspace{0.5cm}

\begin{flushright}
{\small\color{muted} 下一章：请把手弄脏——关于极致野望的七个字。}
\end{flushright}
% ============================================
% 第15章：请把手弄脏——关于极致野望的七个字
% 组合桥 · 终章
% ============================================

\chapter{请把手弄脏}

\vspace{0.5cm}

\begin{center}
{\small\color{muted} 这座桥，来自2025—2026年的课堂、项目、录音与反复修改。}
{\small\color{muted} 所有的公式、所有的故事、所有的失败和骄傲，都指向这五个字。}
\end{center}

\vspace{1cm}

\section{野望}

野望不是``我要考第一''。不是``我要拿奖''。不是``我要比别人强''。

野望是：\textbf{``我要让这座桥站住。''}

2025年10月21日，我在结构设计大赛的赛场上。一个学生做了一个100克的船，载了19公斤。我在课堂上讲这个故事的时候，说``也不代表啥''。但语气里，有一种藏不住的骄傲。不是骄傲排名。是骄傲\textbf{那个东西站住了}。

2026年7月1日，学生在结题汇报里谈到AI画错的弯矩图，也谈到版本覆盖与反复修改。他最后说：``你就真的得是自己会才行，要不然你也不知道AI做得对不对。''这句话没有口号那么整齐，却更接近项目的真相：\textbf{工具可以推进工作，判断不能凭空生成。}

这就是野望。不是超过别人。是\textbf{做成一件自己本来做不成的事}。

\begin{insight}
\label{insight:chap15}

\textbf{野望不是竞争。野望是创造。}

不是``我要比别人强''。是``我要亲手做出一个东西，它以前不存在，现在存在了，而且它站住了''。

这个世界上，最稀缺的不是聪明人。聪明人很多。最稀缺的是\textbf{愿意把手弄脏的人}——愿意花三个小时搭一座桥，塌了五次，第六次站住了，然后抬起头，眼睛里有光，说：``你看，它站住了。''
\end{insight}

\section{技艺的幸福感}

2026年6月，SpanCraft设计复盘。我写下了那八个维度。但八个维度背后，有一个更深的东西——我后来才意识到，那才是SpanCraft真正在做的事。

\textbf{SpanCraft在教学生感受``技艺的幸福感''。}

什么是技艺的幸福感？是你亲手做了一样东西，你知道它每一个细节，你为它花了时间，你为它失败过，你为它骄傲过。你不需要别人告诉你它好不好。你自己知道。你看着它，心里有一种\textbf{安静}的满足。

这种满足，AI给不了你。AI可以生成一座完美的桥，比你搭的漂亮一百倍。但那是AI的桥。不是你的桥。你看着AI的桥，不会心跳加速。你看着\textbf{你自己搭的桥}——那座歪歪扭扭、梁太高、索太粗、但它站住了的桥——你才会心跳加速。

\textbf{因为那座桥上有你的时间。你的错误。你的决定。你的骄傲。}

\section{五个字}

这本书叫《请把手弄脏》。五个字。

这不是一个比喻。我真的是在说：请伸出你的手。去碰混凝土。去弯钢筋。去调参数。去写代码。去搭一座桥。然后看它塌掉。再搭一座。再看它塌掉。再搭一座——直到它站住。

\textbf{不是AI做不了。是AI做了，就不是你的了。}

2026年5月22日，我在游戏化教学分享中说了那句话：``学生在学习过程中从来没有失败的权利，他也就没有成功的那个喜悦或者荣耀。''这句话，被反复引用，被写进了教学理念分析，被印在了这本书的第2章。

但我想在这里再说一遍。因为它是全书最重要的一句话。

\textbf{失败的权利，是学习的权利。}如果你从来没有失败过——如果你从来没有搭过一座会塌的桥，从来没有写过一段会报错的代码，从来没有问过一个``蠢''问题——那你从来没有真正学过任何东西。你只是被灌输过。被灌输的东西，AI可以做得更好。但你自己学到的东西——那些在你失败过、调整过、最终理解了的东西——AI永远拿不走。

\section{桥的本事}

桥的本事，是把原本分开的两岸，连成彼此的两岸。

海德格尔说这句话的时候，他说的不是工程。他说的是\textbf{存在}。一座桥不是简单地连接了A点和B点。一座桥让A点成了``此岸''，让B点成了``彼岸''。在桥出现之前，两岸只是两块土地。在桥出现之后，它们成了彼此的对方。

\textbf{一堂好课，做的也是同一件事。}

在好的课堂出现之前，学生和学生是两个互不相干的个体。在好的课堂出现之后，他们成了``我们''——一群一起搭过桥的人，一起失败过、一起骄傲过、一起在深夜调过参数的人。

\textbf{一个好游戏，做的也是同一件事。}

在好游戏出现之前，``知识''和``我''是两个互不相干的东西。在好游戏出现之后，知识成了我手指上的感觉——那个$qL^2/8$，不再是黑板上的一行公式，而是我亲手让它塌过三次之后，终于理解了的东西。

\textbf{一本好书，做的也是同一件事。}

在好书出现之前，``AI时代的教学''是一个让人焦虑的话题。在好书出现之后——我希望——它变成了一个让人兴奋的邀请。不是``你应该学''。不是``你应该改''。是\textbf{``你想不想来试试？''}

\section{在一个略显荒凉的时代里}

这本书写于2025—2026年。

这不是一个对土木工程特别友好的时代。分数线在掉。新生群里有人劝退。AI可以生成一切``看起来像桥''的东西。

但这也是一个对土木工程特别需要的时代。因为\textbf{当AI能生成一切的时候，真正稀缺的，就不再是生成，而是判断。不再是速度，而是责任。不再是看起来像，而是真的站得住。}

土木工程教会我的，恰好就是这些。

\textbf{你不能骗一座桥。}你可以骗一份PPT。你可以骗一份报告。你可以骗一个评委。你可以骗你自己。但你不能骗一座桥。桥不会听你解释。它只会站住，或者塌掉。

在一个越来越多人相信``差不多就行''的世界里，这种``不能骗''的品格，是一种稀缺的、值得被保护的东西。

\textbf{创造美，并且严肃地为这份美负责。}这大概是工程教育里最难教、也最值得教的东西。而要教会它，第一步，是先让人重新愿意\textbf{把手弄脏}。

\section{带着骄傲，和一点悲伤}

2026年6月，SpanCraft设计复盘。结尾我写了这样一段话：

\begin{shiyu-inline}
``桥的本事，是把分开的两岸连起来。也许一堂好课、一个好游戏，做的也是同一件事。在一个略显荒凉的时代里，把'骄傲'，和那些'还愿意相信的人'，重新连接起来。带着骄傲，和一点悲伤，我把这座桥，建给每一个愿意从它背上走过的人。''

\hfill ——SpanCraft设计复盘，2026年6月
\end{shiyu-inline}

\textbf{带着骄傲，和一点悲伤。}

骄傲，是因为我们做到了。我们证明了：好玩和有料，不矛盾。AI和人的判断，不矛盾。效率和意义，不矛盾。我们证明了：在AI能做一切的时代，教育仍然可以是一件让人心跳加速的事。

悲伤，是因为不是所有人都看到了。还有无数学生在问``学这个有什么用''。还有无数老师在担心``AI会不会取代我''。还有无数课堂，仍然把大量时间用在复述那些工具已经可以迅速生成的内容上。

但这本书，就是一座桥。

让一个从没想过土木的人，因为它好看好玩，愿意多停留三十秒。让一个正在犹豫要不要转专业的学生，忽然想起来——结构，原来是有美的。让一个被AI吓到的教师，重新发现——\textbf{工具能力在变，但人仍要决定什么值得教、如何核验，以及由谁负责}。

\begin{center}
{\LARGE\bfseries 请把手弄脏。}
\end{center}

\vspace{2cm}

\begin{flushright}
ColdCat（大连交通大学）\\
2026年7月，于大连
\end{flushright}


% ====== 尾声 ======
\chapter*{尾声：把下一次修改留下来}
\addcontentsline{toc}{chapter}{尾声：把下一次修改留下来}

这本书最初想用一个漂亮故事收尾：一个学生写出有限元求解器，发现三个错误，最后说出一句关于“让桥站住”的话。故事完整、动人，也恰好证明全书的每一个观点。

可惜，在现有录音转写里找不到它。

如果一本讨论 AI 时代专业判断的书，需要靠一段无法核实的合成故事完成高潮，它已经背叛了自己的主题。因此，那个故事被删掉了。留下来的现场没有那么整齐，却更值得写。

2026 年 7 月 1 日的结题汇报持续近三个小时。学生展示了完成度差异很大的作品：有的模型能运行，却解释不清公式；有的游戏有新想法，但自动搭桥始终做不出来；有的小组发现 AI 把弯矩图画错，于是回到教材和课堂内容，逐个指定关键点；也有人坦率地说，继续迭代时把较好的半成品覆盖了。

这不是“AI 赋能教学”的整齐样板。它更像真实工程：工具有用，版本会坏，模型会错，时间有限，人必须作判断。

其中一个学生总结：

\begin{shiyu-inline}
``你就真的得是自己会才行，要不然你也不知道 AI 做得对不对。''

\hfill ——课程结题汇报，2026 年 7 月 1 日 00:25:20
\end{shiyu-inline}

这句话并不意味着学生应该先把所有知识学完，才有资格使用 AI。相反，项目把“还不会什么”暴露了出来。问题是课堂是否给学生第二次机会：回到来源，修改作品，解释变化，并把错误变成下一次判断的依据。

反馈研究把这一步称为关闭反馈回路：信息只有被学习者理解并用于后续行动，才成为反馈\cite{carless2019spirals}。这也解释了为什么本书反复强调修改记录。最终成品容易伪装，修改过程更接近学习发生的地方。

如果你是一位教师，读到这里不必立刻开发游戏、网站或 AI 代理。请先保存一份学生的初稿，再保存修改后的版本；问一句“你改了什么，依据是什么”；把回答写进下一次任务的评价标准。

如果你是一位学生，也请不要只保存最好看的结果。保留那张画错的图、那个无法运行的版本和你后来查到的依据。它们不是羞耻的残骸，而是专业判断长出来的年轮。

\enlargethispage{2\baselineskip}
桥梁工程里，验收不靠一张渲染图。学习也不该只靠一个最后答案。我们真正要留下的，是从第一次尝试到下一次修改之间那段可追溯的路。下一节课，就从那里开始。


% ====== 附录 ======
\appendix
\chapter{一页式学习回路设计表}

这张表用于重做一个 45--90 分钟教学单元。先填写“学生要做什么”，最后才填写“教师讲什么”。

\begin{description}[style=nextline,leftmargin=3.2cm]
  \item[值得求解的问题] 学生为什么愿意解决？问题有哪些真实约束？
  \item[目标判断] 课后学生应能作出什么专业判断，而不仅是记住什么？
  \item[首次作品] 15 分钟内能产出什么可见结果？允许粗糙，但必须可检查。
  \item[反馈来源] 结构响应、同伴、量规、教师或 AI 分别反馈什么？
  \item[第二次修改] 学生必须根据反馈改变哪一处，并保存前后版本。
  \item[解释与答辩] 学生怎样说明“改了什么、依据是什么、仍不确定什么”？
  \item[AI 边界] AI 可以生成哪些材料？哪些判断必须由学生亲自完成？
  \item[验收证据] 初稿、反馈、修改稿、来源和口头解释是否齐全？
\end{description}

\section*{最小示例：简支梁跨中弯矩}

不要先让学生抄写公式。先给三座跨度不同、截面相同的简支梁，让学生预测哪一座最危险并画出弯矩图形状；再用计算或交互工具反馈；最后要求学生用 $qL^2/8$ 解释跨度为何如此敏感。AI 可以帮助画图和检查代数，但学生必须说明边界条件、单位和图形特征。

\section*{验收标准}

一个合格单元至少包含一次真实修改。如果学生只生成一次答案、教师只给一次分数，这不是学习回路，只是一次提交。

\chapter{AI 参与工程作业的证据链}

AI 使用声明不应只问“用了没有”，而应回答“在哪一步用了、输入了什么、怎样核验、谁承担结论”。建议学生随作业提交以下六项。

\begin{enumerate}
  \item \textbf{任务定义：}问题、约束、预期输出和评价标准；
  \item \textbf{来源：}教材、规范、数据与参考模型的版本；
  \item \textbf{AI 记录：}关键提示、工具名称、日期和必要的输出节选；
  \item \textbf{核验：}独立计算、极限情况、量纲、对照案例或第二工具复核；
  \item \textbf{修改：}AI 初稿与最终稿之间最重要的三处变化；
  \item \textbf{责任声明：}仍然存在的不确定性，以及最终判断由谁作出。
\end{enumerate}

\section*{四级量规}

\begin{description}[style=nextline,leftmargin=2.2cm]
  \item[不足] 只有最终成品；来源、AI 参与和核验过程不可追溯。
  \item[基本] 声明工具并保留部分过程，但核验主要停留在“能运行”。
  \item[可靠] 关键假设、单位、边界条件和结果均有独立检查；能指出 AI 错误。
  \item[成熟] 比较多种方案，解释取舍与不确定性，并能在追问中重建判断过程。
\end{description}

量规评的是证据质量，不是提示词长度。提示词写得很长，不代表工程判断可靠。

\chapter{课堂录音进入公开写作的处理流程}

课堂原声能让教学写作摆脱空泛，但也带来转写误差、隐私与选择性引用风险。建议采用以下流程。

\begin{enumerate}
  \item \textbf{归档：}保留原始音频、转写文件、日期、时长和来源平台；
  \item \textbf{定位：}每条候选引文记录文件名和时间戳；
  \item \textbf{回听：}正式出版前回听关键句，修正自动转写错误；
  \item \textbf{去识别：}删除姓名、班级、账号、个人经历等非必要信息；
  \item \textbf{取授权：}公开学生作品、长篇发言或可识别经历前取得授权；
  \item \textbf{保语境：}同时记录前后问题，避免把犹豫、玩笑和现场估计写成定论；
  \item \textbf{分证据：}教师分享属于实践反思，学生发言属于个案材料，均不自动成为效果数据；
  \item \textbf{建台账：}对每条公开引文记录状态：待回听、已回听、已授权、仅内部使用。
\end{enumerate}

\begin{tech-note}
公开 GitHub 仓库不应上传原始课堂录音和未经匿名化的完整转写。书中可保留日期与时间戳，原始证据放在受控的私有档案中。
\end{tech-note}

\chapter{教学矛盾决策表}

本表借用 TRIZ 的“先说清矛盾，再寻找分离或转化条件”思路。它是备课启发工具，不是经过验证的教学理论。

\section*{开放探索 / 专业正确}
常见失败是完全放任或直接给答案。可以先让学生自由生成方案，再设置强制核验关。检查：学生何时必须回到规范、量纲或边界条件？

\section*{即时反馈 / 独立思考}
常见失败是 AI 过早给出答案。可以先提交预测，再开放提示与计算。检查：学生是否在看答案前作过一次可见尝试？

\section*{个性化 / 教师负担}
常见失败是无限制逐份精批。可以让工具检查低层错误，教师聚焦关键判断。检查：教师反馈是否用在只有人能承担的地方？

\section*{游戏吸引 / 学习目标}
常见失败是积分热闹、知识游离。可以让核心操作本身产生学科反馈。检查：得高分是否真的需要使用目标知识？

\section*{作品精美 / 过程真实}
常见失败是最终图像遮蔽错误。可以强制提交初稿、版本差异和口头答辩。检查：能否看见作品如何变成现在的样子？

使用时只选一对矛盾。先写清当前课程在哪一端失衡，再设计一个最小改动，并用学生作品验证，而不是一次引入所有新工具。

\chapter{SpanCraft 原型复盘：从概念堆积到单一回路}

SpanCraft 是一个桥梁建造教学原型。它的价值不在于已经证明了学习效果，而在于暴露了教学软件常见的设计错误。

\section*{第一版的问题}

第一版试图同时呈现材料、截面、成本、受力、规范和桥型。信息都“有用”，但玩家打开界面先面对参数和说明，尚未形成一个想解决的问题。结果更像把作业搬进网页。

\section*{重排原则}

修订时采用两个原则：先给魅力，再给负荷；一次只突出一个可感知的关系。玩家先搭、先试车、先看见变形或破坏，再逐步接触公式和约束。可视化放大只改变显示，不应伪装成真实变形；所有演示性处理都需要标明。

\section*{尚未完成的验证}

当前原型能够说明一种设计思路，但尚不能回答“是否提高了学习成绩”“效果能否迁移”“不同基础学生是否同样受益”。后续验证至少需要：明确学习目标、前后测或对照任务、过程日志、访谈，以及对错误概念的检查。

\section*{可复用的产品检查表}

\begin{itemize}
  \item 玩家每次操作会收到什么学科反馈？
  \item 不懂目标知识，是否也能靠反复点击过关？
  \item 显示放大、简化模型和真实计算是否清楚区分？
  \item 失败是否告诉玩家下一步可以检查哪里？
  \item 作品能否导出，供学生解释与教师评价？
\end{itemize}

\chapter{资源索引与版本说明}

\section*{公开资源}

\begin{sloppypar}
\begin{itemize}
  \item SpanCraft 桥造计划：\url{https://syzhaoln-stack.github.io/spancraft-bridge-game/}
  \item 本书仓库：\url{https://github.com/syzhaoln-stack/build-a-bridge}
\end{itemize}
\end{sloppypar}

\section*{本地参考资产}

作者的私有工作区还包含桥梁数学可视化、卡牌游戏设计、桥梁网站、课程备课、抗风课程材料和教师培训案例。它们不会整体公开进入本仓库。正文采用其中内容时，将注明原项目、版本日期、改动范围和公开授权状态。

\section*{引用规则}

网页和 AI 工具变化很快。复现案例时请同时记录访问日期、模型或软件版本、关键配置和本地存档。书中原则性方法不绑定某个产品；只有在复现需要时才保留产品名。

\section*{勘误与贡献}

欢迎通过 GitHub Issue 提交以下内容：可定位的事实错误、无法复现的步骤、工程概念异议、课堂迁移案例。纯粹的工具推广和没有证据的效果宣称不进入正文。


% ====== 参考文献 ======
\printbibliography[title={参考文献}]

\end{document}
